\subsection{The positive-degree comparison and its fiber over global fields}\label{sec:global-spectrum}
In this subsection, we construct the comparison $g\colon\tau_{\ge1}K(D)\to\tau_{\ge1}K(F)$ over a global field (Theorem~\ref{thm:paper-global-comparison}) and study its fiber $Z$.
We split $g$ in degrees at least two, first over number fields and then over global function fields, and compute $K_3(D)$.
We then show that $Z$ has finitely many primary factors with bounded homotopy (Theorem~\ref{thm:paper-global-fiber-bounds}), that, up to equivalence, it depends only on the cyclic subgroup of the Brauer group generated by $[D]$ (Proposition~\ref{prop:paper-brauer-power-fiber}), and that over number fields it is noncanonically equivalent to a finite sum of shifted Eilenberg--Mac Lane spectra (Theorem~\ref{thm:paper-global-fiber-gem}).
When $D$ is ramified at a real place, the statements about the comparison, its splitting, and its fiber are made after inverting $2$; the computation of $K_3(D)$ is integral.

\subsubsection{Construction of the positive-degree comparison}
The comparison is glued in the arithmetic square from the inverse of multiplication by a class of geometric rank $1$ at each prime $\ell\ne\operatorname{char}F$ dividing $\deg D$ (Proposition~\ref{prop:global-rank}), from $\Nrd_1$ at $p=\operatorname{char}F$ when $p\mid\deg D$, and from $N_{D/F}/\deg D$ after inverting $\deg D$.

Let $F$ be a global field.
In the construction below we assume that $D$ splits at all real places; the real-ramified case is treated after inverting $2$ in Theorem~\ref{thm:paper-global-comparison}.
As in Proposition~\ref{prop:global-real-weight-two}, $H^i_{\mathrm{rel}}$ denotes the cohomology of the fiber of restriction to the algebraic real closures at all real places.
Fix a prime $\ell\ne\operatorname{char}F$.
In \'etale descent, the geometric stalks with coefficients modulo $\ell^r$ are
\[
 \Z/\ell^r(n)\ \text{in degree }2n,\qquad 0\ \text{in odd degrees}.
\]
Given a compatible geometric rank-$1$ class $x\in\pi_0L_{K(1),\ell}K(D)$, multiplication by $x$ is an equivalence
\[
 m_x\colon L_{K(1),\ell}K(F)\xrightarrow{\ \sim\ }L_{K(1),\ell}K(D)
\]
by Lemma~\ref{lem:paper-etale-rank-one} (see \cite[Theorems 6.18 and 7.14]{ClausenMathew}).

\begin{proposition}\label{prop:global-rank}
Let $F$ be a global field, let $\ell\ne\operatorname{char}F$ be a prime, and assume that $D$ splits at all real places of $F$.
Then one can choose a geometric rank-$1$ class $x\in\pi_0L_{K(1),\ell}K(D)$.
For a number field with real places and for $\ell=2$, one can prescribe the rank-$1$ Morita class at every real closure.
\end{proposition}
\begin{proof}
We have $\cd_\ell F\le2$ unless $\ell=2$ and $F$ has a real place; this exceptional case is the real-relative case treated below.

Suppose first that $\cd_\ell F\le2$.
Then the rank lifts along the edge map in total degree zero modulo $\ell^r$, and the difference of two lifts lies in $H^2(F,\Z/\ell^r(1))$.
The coefficient transition maps
\[
 H^2(F,\Z/\ell^{r+1}(1))\longrightarrow H^2(F,\Z/\ell^r(1))
\]
are surjective, so we can choose compatible lifts.

In the real-relative case, the obstruction to lifting the prescribed real elements lies in total degree $-1$ of the relative groups, namely in $H^1_{\mathrm{rel}}$ in weight zero.
It is the relative boundary of the tuple
\[
 (1)_v\in\prod_{v\ \mathrm{real}}H^0(F_v^{\mathrm{rc}},\Z/2^r),
\]
which is the restriction of the diagonal $1$, so this boundary vanishes.
The differences of lifts lie in relative $H^2$ in weight one, whose transition maps are surjective because relative $H^3$ vanishes.
The Milnor exact sequence for the homotopy inverse limit therefore supplies $x$.
Since the corresponding groups on the real side modulo $2^r$ are finite, the restriction of $x$ is still the prescribed one in the limit.
\end{proof}

We also write $N_{D/F}\colon L_{K(1),\ell}K(D)\to L_{K(1),\ell}K(F)$ for the map induced by the norm, and $m_{N_{D/F}(x)}$ for multiplication by $N_{D/F}(x)\in\pi_0L_{K(1),\ell}K(F)$.

\begin{lemma}\label{lem:paper-global-product-vanishing}
Let $\ell\ne\operatorname{char}F$, and let $x\in\pi_0L_{K(1),\ell}K(D)$ be a class of geometric rank $1$ as in Proposition~\ref{prop:global-rank}; if $\ell=2$ and $F$ has real places, we assume that $x$ is the rank-$1$ Morita class at every real closure.
Then
\[
 \bigl(N_{D/F}(x)-\deg D\bigr)\,y=0\qquad\text{for }y\in\pi_iL_{K(1),\ell}K(F),\ i\ge1,
\]
so $N_{D/F}\circ m_x$ acts as multiplication by $\deg D$ on $\pi_i$ for $i\ge1$.
Moreover, for every finite separable splitting field $E/F$, multiplication by $x_E-1$ vanishes on $\pi_iL_{K(1),\ell}K(E)$ for $i\ge1$, where $x_E=\res_{E/F}x$ and $K(D_E)$ is identified with $K(E)$ by Morita equivalence.
\end{lemma}
\begin{proof}
By the projection formula,
\[
 N_{D/F}\circ m_x\simeq m_{N_{D/F}(x)}.
\]
The classes $N_{D/F}(x)-\deg D\in\pi_0L_{K(1),\ell}K(F)$ and $x_E-1\in\pi_0L_{K(1),\ell}K(E)$ have geometric rank $0$.
In the real-relative case they restrict to zero at every real closure, because $x$, and hence $x_E$, is the rank-$1$ Morita class there.
Lemma~\ref{lem:paper-rank-zero-product}, applied over $F$ and over $E$, gives both assertions.
\end{proof}

\begin{lemma}\label{lem:paper-global-no-divisible}
Let $F$ be a global field and $j\ge1$.
Then $K_j(F)$ has no nonzero divisible subgroup.
The same holds for $K_j(F)_{(\ell)}$ for every prime $\ell$ and for $K_j(F)[1/2]$.
\end{lemma}
\begin{proof}
In degree one, we use the map to the direct sum of the valuations at the finite places or closed points; its kernel is the finitely generated unit group.
In higher degrees, we use localization for the ring of integers or the smooth proper curve: the direct sum of the groups $K_{j-1}$ of the residue fields has no nonzero divisible subgroup, so any divisible subgroup lies in the image of the finitely generated group $K_j$ of the model \cite{QuillenHigher,QuillenCurve,QuillenFinite}, \cite[Theorem 1]{QuillenFiniteGeneration}.
After localization at $\ell$ or after inverting $2$, the direct sum of the value groups or of the residue $K$-groups still has no nonzero divisible subgroup, and the rest lies in a finitely generated $\Z_{(\ell)}$-module or $\Z[1/2]$-module, so the same argument applies.
\end{proof}

\begin{theorem}\label{thm:paper-global-comparison}
If $D$ splits at all real places of $F$, there is an integral map of positive-degree spectra
\[
 g\colon\tau_{\ge1}K(D)\longrightarrow\tau_{\ge1}K(F)
\]
that carries the splitting transfer $\theta_E$ from each finite separable splitting field $E/F$ to $N_{E/F}$ on every $\pi_j$.
In degrees $1$ and $2$, the map $g$ agrees with the reduced norms $\Nrd_1$ and $\Nrd_2$.
With real ramification the same construction works after inverting $2$.
\end{theorem}

For $j\ge1$, we write $g_j=\pi_jg$.

\begin{proof}
We first define the components of $g$ in the arithmetic square.
For $\ell\mid\deg D$ with $\ell\ne\operatorname{char}F$, we use the composite
\[
 K(D)\longrightarrow L_{K(1),\ell}K(D)\xrightarrow{\ m_x^{-1}\ }L_{K(1),\ell}K(F),
\]
together with the Quillen--Lichtenbaum comparison $\tau_{\ge1}L_{K(1),\ell}K(F)\simeq(\tau_{\ge1}K(F))^\wedge_\ell$ for $F$.
If $p=\operatorname{char}F$ divides $\deg D$, then $[F:F^p]=p$, and the rank-$1$ Frobenius comparison of Proposition~\ref{prop:app-frobenius} gives
\[
 K_i(D;\Z/p)=K_i(F;\Z/p)=0\qquad(i\ge2).
\]
Thus the completed positive-degree spectrum at $p$ is concentrated in degree one, where we use the global reduced norm $\Nrd_1$.
After inverting $\deg D$, we use $N_{D/F}/\deg D$.

These components agree on the overlap.
For $\ell\ne\operatorname{char}F$, Lemma~\ref{lem:paper-global-product-vanishing} gives $N_{D/F}m_x=\deg D$ in positive degrees, so
\[
 m_x^{-1}=N_{D/F}/\deg D
\]
on the rationalized positive homotopy groups of the $\ell$-completed factor.
If $p=\operatorname{char}F$ divides $\deg D$, then for every finite separable splitting field $E/F$
\[
 N_{D/F}\theta_E=\deg D\cdot N_{E/F}=\deg D\cdot\Nrd_1\theta_E,
\]
and the splitting transfers generate $K_1(D)$, as shown below, so $N_{D/F}=\deg D\cdot\Nrd_1$ on $K_1(D)$.
Since the overlap of the arithmetic square is rational, maps into it that agree on homotopy groups are homotopic, so we can glue the completed maps and $N_{D/F}/\deg D$ to obtain $g$.
With real ramification, after inverting $2$ only odd primes $\ell$ occur, and then $\cd_\ell F\le2$.
For odd $\ell$, the rank-$1$ lift, the multiplicative equivalence and the vanishing of the products in Lemma~\ref{lem:paper-global-product-vanishing} follow by the same argument from absolute \'etale descent; the geometric rank argument uses no prescription at the real places.

Next we check the transfer identity on the components.
For a finite separable splitting field $E$, we identify $K(D_E)$ with $K(E)$ by the Morita equivalence in the definition of $\theta_E$, so that $\theta_E=N_{D_E/D}$.
On the completed factors, the projection formula of Lemma~\ref{lem:paper-etale-rank-one} gives $N_{D_E/D}m_{x_E}\simeq m_xN_{E/F}$, and hence
\[
 m_x^{-1}\theta_E=N_{E/F}m_{x_E}^{-1}.
\]
By Lemma~\ref{lem:paper-global-product-vanishing}, multiplication by $x_E-1$ vanishes in positive degrees over $E$, so the right-hand side induces the field norm.
For the only positive degree of the completed characteristic-$p$ factor, the transfer identity for $\Nrd_1$ gives the same conclusion.
After inverting $\deg D$, we use $N_{D/F}\theta_E=\deg D\cdot N_{E/F}$.
Consequently, on integral homotopy groups $g_j\theta_E-N_{E/F}$ takes values in the divisible image of the boundary from the rational overlap of the arithmetic square.
This image vanishes by Lemma~\ref{lem:paper-global-no-divisible}, applied to $K_j(F)$ or, with real ramification, to $K_j(F)[1/2]$, so the transfer identity holds with the coefficients of the theorem.

Finally we identify $g_1$ and $g_2$.
In degree one, the units generating separable maximal subfields form an open subset, and every unit is the ratio of two elements of this subset; consequently the splitting transfers generate $K_1(D)$.
In degree two, the proof of Theorem~\ref{thm:paper-global-norm} shows that $K_2(D)$ is generated by transfer from one separable splitting field.
Since the reduced norms satisfy the same transfer identities $\Nrd_j\theta_E=N_{E/F}$ for $j=1,2$, we obtain $g_j=\Nrd_j$.
\end{proof}

\begin{lemma}\label{lem:paper-global-low-degrees}
Let $F$ be a global field, and let $g$ be the comparison of Theorem~\ref{thm:paper-global-comparison}.
If $D$ splits at all real places of $F$, then $g_1$ and $g_2$ are isomorphisms; otherwise they are isomorphisms after inverting $2$.
\end{lemma}
\begin{proof}
By Theorem~\ref{thm:paper-global-comparison}, $g_j=\Nrd_j$ for $j=1,2$.
Over a global function field, $\SK_1(D)=0$ by \cite[Corollary 2]{Yanchevskii1975}, whose proof also gives surjectivity of $\Nrd_1$.
Over a number field, $\Nrd_1$ is injective because $\SK_1(D)=0$ \cite{Wang1950}.
By the Hasse--Schilling--Maass norm theorem \cite[Theorem 33.15]{ReinerMaximalOrders}, its image consists of the elements of $F^\times$ that are positive at every real place where $D$ is ramified.
Hence $\Nrd_1$ is surjective if $D$ splits at all real places.
If $D$ is ramified at a real place, the cokernel of $\Nrd_1$ is killed by $2$, since squares are positive at every real place, so $\Nrd_1$ is an isomorphism only after inverting $2$.
By Theorem~\ref{thm:paper-global-norm}, $\Nrd_2$ is injective and its cokernel is $(\Z/2)^{S_{\mathbb R}}$, where $S_{\mathbb R}$ is the set of real places where $D$ is ramified.
Thus $\Nrd_2$ is an isomorphism if $D$ splits at all real places, and after inverting $2$ in general.
\end{proof}

Over number fields, the comparisons of Theorem~\ref{thm:paper-global-comparison} can be chosen coherently over the intermediate fields of a finite Galois extension, as in the local case (Proposition~\ref{prop:paper-finite-galois-naturality}).

\begin{proposition}\label{prop:paper-finite-galois-naturality-global}
Let $F$ be a number field, let $D$ be a central division algebra over $F$, and let $L/F$ be a finite Galois extension.
If $D$ is ramified at some real place, we invert $2$ throughout.
Then the comparisons
\[
 g_{F'}\colon\tau_{\ge1}K(D_{F'})\longrightarrow\tau_{\ge1}K(F')
\]
of Theorem~\ref{thm:paper-global-comparison} for the intermediate fields $F'$ of $L/F$ can be chosen simultaneously and coherently with respect to restriction, transfer, conjugation, and their composites.
\end{proposition}
\begin{proof}
Let $G=\Gal(L/F)$, and define $X,Y\in\mathscr M_G$ by $X(G/H)=\tau_{\ge1}K(D_{L^H})$ and $Y(G/H)=\tau_{\ge1}K(L^H)$, as in the proof of Proposition~\ref{prop:paper-finite-galois-naturality}.
As there, we use the norms $N_{D_{F'}/F'}$ of the algebras $D_{F'}$, which have degree $\deg D$ even when the index drops over $F'$, and the arithmetic square in $\mathscr M_G$ for the primes dividing $\deg D$.
Only if $D$ is ramified at some real place do we replace $X$ and $Y$ by their localizations away from two at the outset; then the two-completed factor is zero, and we omit it in choosing rank-$1$ classes.

For each remaining $\ell\mid\deg D$, choose a class $x_\ell$ of geometric rank $1$ over $F$ and restrict it to every intermediate field.
If $D$ splits at all real places, the class is supplied by Proposition~\ref{prop:global-rank}, and Lemma~\ref{lem:paper-etale-rank-one} applies; if moreover $\ell=2$, we choose $x_2$ with the rank-$1$ Morita class prescribed at every real closure, as Proposition~\ref{prop:global-rank} allows.
This prescription is inherited by the real closures of the intermediate fields and of the splitting fields.
Under real ramification only odd $\ell$ remain; there we do not apply Proposition~\ref{prop:global-rank} or Lemma~\ref{lem:paper-etale-rank-one} directly, which assume that $D$ splits at the real places, but lift the rank-$1$ class by absolute \'etale descent as in the proof of Theorem~\ref{thm:paper-global-comparison}, and construct the equivalence by the geometric-point and free-module arguments in the proof of Lemma~\ref{lem:paper-etale-rank-one}.
On each orbit, $u_\ell$ (as in the proof of Theorem~\ref{thm:paper-local-spectrum}) is the composite of the natural localization unit on the twisted side, $m_{x_\ell}^{-1}$, and the inverse of the Quillen--Lichtenbaum comparison on the field side; the latter is natural and an equivalence on each orbit, so it can be inverted in the Mackey category.
The gluing then follows the argument of Proposition~\ref{prop:paper-finite-galois-naturality}.
Here the product calculation in the proof of Theorem~\ref{thm:paper-global-comparison} identifies the rationalized $u_\ell$ with $N/\deg D$ on positive homotopy groups over every intermediate field.
We obtain $g\colon X\to Y$ from the pullback.

For number fields the correction of the local case is unnecessary.
Once the arithmetic components are fixed, the maps on homotopy groups in degrees $j\ge1$ are determined up to the boundary image of the rational group $\pi_{j+1}\bigl(Y(F')^\wedge_{\deg D}[1/\deg D]\bigr)$, and this image is divisible and hence zero by Lemma~\ref{lem:paper-global-no-divisible}, applied to $K_j(F')$ or, under real ramification, to $K_j(F')[1/2]$.
For each intermediate field $F'$ and each finite separable splitting field $E/F'$, the projection formula and the vanishing of products with classes of geometric rank $0$, as in the proof of Theorem~\ref{thm:paper-global-comparison}, give the transfer identity on the completed components.
After inverting $\deg D$, the identity follows from $\pi_j(N_{D_{F'}/F'}\circ\theta_E)=\deg D\cdot\pi_jN_{E/F'}$.
The vanishing of the boundary image therefore gives the transfer identity on homotopy groups with the stated coefficients.
Applying the generation by splitting transfers, as in the proof of Theorem~\ref{thm:paper-global-comparison}, to the division algebra representing $D_{F'}$ over each intermediate field gives the degree-one and degree-two normalizations.
Hence $g_{F'}$ satisfies the transfer identities and the degree-one and degree-two normalizations of Theorem~\ref{thm:paper-global-comparison}.
\end{proof}

\subsubsection{Derived norms and splitting over number fields}
Over number fields, the splitting of Theorem~\ref{thm:paper-global-split} uses the derived norm equivalence of Proposition~\ref{prop:paper-derived-norm} for a finite Galois extension $E/F$, with every real place of $F$ split completely in $E$ unless $2$ is inverted.

For a finite Galois extension $E/F$ with group $G$, the field norm defines a map from homotopy orbits using the canonical isomorphism between restriction of scalars and its composite with Galois conjugation.
We include the coherence of these isomorphisms in the action defining the orbits.

We use \'etale cohomology of rings of $S$-integers, as follows.
Let $E/F$ be a finite Galois extension of number fields with group $G$.
Fix $\ell$, take a finite set $S$ of places containing the places over $\ell$ and those ramified in $E/F$, and put
\[
 V_L=\Spec\mathcal O_{L,S}\qquad(L\in\{E,F\}).
\]
Then $V_E\to V_F$ is a finite \'etale $G$-cover.
For a finite $\ell$-primary local system $M$, we use ordinary $R\Gamma$ if $\ell$ is odd or $F$ is totally imaginary.
If $\ell=2$ and $F$ has real places, we use
\[
 R\Gamma_{\mathrm{rel}}(V_L,M)=
 \fib\!\left(R\Gamma(V_L,M)\longrightarrow
              \bigoplus_{v\ \mathrm{real}}R\Gamma(L_v,M)\right).
\]
The real terms are ordinary group cohomology of $\Z/2$.
We write $R\Gamma_\star(V_L,-)$ for the chosen functor.
The cohomology of $R\Gamma_\star(V_L,M)$ lies in $[0,2]$ for finite $\ell$-primary local systems because $S$ contains the places above $\ell$.
For $R\Gamma_{\mathrm{rel}}(V_L,M)$, this bound uses the isomorphism of real restriction in degrees at least three and surjectivity in degree two.
Surjectivity in degree two follows from Poitou--Tate duality: choose a place $v\mid\ell$ of $L$; local duality makes $H^2(L_v,M)\to H^0(V_L,M^D)^\vee$ surjective, where $M^D=\operatorname{Hom}(M,\mu_{\ell^\infty})$, so the obstruction for any tuple at the real places can be cancelled at that finite place.
For continuous $\Z_\ell$ coefficients, finiteness at each finite stage gives the Mittag--Leffler condition and preserves this bound \cite{KahnPositive,ChinburgEtal}.

\begin{lemma}\label{lem:paper-derived-tate-norm}
Let $M$ be a finite local system of $\Z/\ell^a$-modules on $V_F$, and let $\omega\colon V_E\to V_F$ be the cover.
In the real-relative case, assume that every real place of $F$ splits completely in $E$.
Then the trace induces an equivalence
\begin{equation}\label{eq:paper-derived-tate-norm}
  \Z/\ell^a\otimes_{\Z/\ell^a[G]}^{\mathbf L} R\Gamma_\star(V_E,\omega^*M)
      \xrightarrow{\ \sim\ } R\Gamma_\star(V_F,M).
\end{equation}
\end{lemma}
\begin{proof}
For the cover $\omega$, the stalks of $\omega_*\omega^*M$ are induced $G$-modules with derived coinvariants $M$ via trace.
This is the tensor product calculation with the free group-ring resolution for a regular permutation action.
We have $R\Gamma(V_F,\omega_*\omega^*M)=R\Gamma(V_E,\omega^*M)$, compatibly with restriction to real places in the real-relative case because they split completely.
A free resolution then gives~\eqref{eq:paper-derived-tate-norm}: finite truncations of the bar resolution have finite free terms, so sections commute with the finite sums and cofibers.
The omitted tail contributes nothing to a fixed degree past a sufficiently late stage by the cohomological upper bound two.
The same reasoning applies to relative sections.
\end{proof}

\begin{proposition}\label{prop:paper-derived-norm}
For a finite Galois extension $E/F$ of number fields in which every real place splits completely, the norm
\[
  \mathcal N_{E/F}:\bigl(\tau_{\ge2}K(E)\bigr)_{hG}\longrightarrow \tau_{\ge2}K(F)
\]
is an equivalence.
Without complete splitting at the real places the same holds after inverting $2$.
\end{proposition}
\begin{proof}
Fix $\ell$ and $S$ as above.
The coefficient triangles comparing continuous Tate coefficients with reduction modulo $\ell^a$ are compatible with sections.
The equivalences of Lemma~\ref{lem:paper-derived-tate-norm} modulo $\ell^a$ therefore give the corresponding equivalence in the derived $\ell$-complete category.

We now pass to \'etale $K$-theory.
Write $\mathcal{KU}_\ell$ for the periodic \'etale $K$-theory sheaf, and consider the geometric Postnikov tower
\[
  \mathcal P_L^m=R\Gamma_\star(V_L,\tau_{\ge2m}\mathcal{KU}_\ell),
  \qquad
  \gr^m\mathcal P_L \simeq R\Gamma_\star(V_L,\Z_\ell(m))[2m].
\]
We interpret the continuous layers by inverse limits of their finite-coefficient versions.
The connectivity bound for $\mathcal P_L^m$ is $2m-2$.
Since $H^0(V_L,\Z_\ell(1))=0$ and relative $H^0$ vanishes as well, weights at most one make no contribution in degrees at least two.
Hence $\mathcal P_L^2\simeq\tau_{\ge2}R\Gamma_\star(V_L,\mathcal{KU}_\ell)$.
Homotopy orbits preserve the lower connectivity bound of the tails, so arbitrarily high layers cannot contribute in any fixed degree \cite{ClausenMathew}.
Applying Lemma~\ref{lem:paper-derived-tate-norm} to the finite-coefficient versions of the layers and inducting on finite truncations, we obtain $(\mathcal P_E^2)_{hG}\simeq\mathcal P_F^2$ in the derived $\ell$-complete category.

Next we return to algebraic $K$-theory.
For ordinary sections, we use the $\ell$-complete comparison $K(V_L)\to K^{\et}(V_L)$ in degrees at least two.
In the real-relative case, we restore the real terms.
After $S$ is enlarged, weak approximation supplies $S$-units $u$ realizing all real sign vectors.
The symbols $\{-1,u\}\in K_2(V_L)$ hit the generators of $\pi_2(K(L_v)^\wedge_2)\simeq\pi_2(KO^\wedge_2)=\Z/2$ at the specified real places, so the restriction $K(V_L)^\wedge_2\to\bigoplus_{v\ \mathrm{real}}KO^\wedge_2$ is surjective on $\pi_2$.
In the commutative square formed by this restriction, the restriction $R\Gamma(V_L,\mathcal{KU}_2)\to\bigoplus_{v\ \mathrm{real}}KO^\wedge_2$, and the comparison maps, the vertical maps are isomorphisms in degrees at least two, so the lower restriction is also surjective on $\pi_2$.
Hence the relative, absolute, and real terms still form fiber sequences after taking $2$-connective covers, and the algebraic relative term is identified with $\tau_{\ge2}R\Gamma_{\mathrm{rel}}(V_L,\mathcal{KU}_2)$ \cite[Theorems 7.12--7.13]{ClausenMathew}.
On $V_E$ the real terms are a finite sum of induced $G$-spectra and the derived norm is an equivalence on them.
The equivalence on the relative terms then gives it on the absolute terms \cite{RognesWeibel,ClausenMathew}.

Finally, we reduce these equivalences modulo $\ell$ and pass to the filtered colimit over $S$.
The $K$-theory, connective covers, orbits, and reduction modulo $\ell$ all commute with this colimit, giving an equivalence over the fields modulo every $\ell$.
Rationally the norm is an equivalence by the restriction-transfer identities and averaging.
The integral equivalence follows from the vanishing of the cofiber rationally and modulo each prime.
After two is inverted, only the odd primes and the rational comparison are needed.
\end{proof}

We next construct the transfer from homotopy orbits that appears in the splitting theorems.
Let $E/F$ be a finite Galois splitting field of $D$ with group $G$.
The group $G$ acts on the category of right $D_E$-modules by semilinear scalar change fixing $D$.
Restriction of scalars from $D_E$ to $D$ has natural isomorphisms with its composites with this action.
These isomorphisms preserve the underlying $D$-module and are coherent with group multiplication, so they give the transfer
\[
 \mathcal N^{D}_{E/F}:\bigl(\tau_{\ge2}K(D_E)\bigr)_{hG}\longrightarrow \tau_{\ge2}K(D)
\]
from homotopy orbits.
For $j\ge2$, its composite with the edge map $K_j(D_E)\to\pi_j\bigl((\tau_{\ge2}K(D_E))_{hG}\bigr)$ is restriction of scalars, that is, the splitting transfer $\theta_E$ under the identification $K(D_E)\simeq K(E)$ by Morita equivalence.

In both splitting theorems, we obtain a section of $\tau_{\ge2}g$ once its composite with the transfer from homotopy orbits is an equivalence.

\begin{lemma}\label{lem:paper-split-from-norm}
Let $F$ be a global field, let $g$ be the comparison of Theorem~\ref{thm:paper-global-comparison}, put $Z=\fib(g)$, and let $E/F$ be a finite Galois splitting field of $D$ with group $G$.
If $D$ is ramified at a real place, invert $2$ throughout.
Assume that $h=(\tau_{\ge2}g)\circ\mathcal N^{D}_{E/F}$ is an equivalence.
Then $\mathcal N^{D}_{E/F}\circ h^{-1}$ is a section of $\tau_{\ge2}g$, so that
\[
 \tau_{\ge2}K(D)\simeq\tau_{\ge2}K(F)\oplus\fib(\tau_{\ge2}g)
 \quad\text{with first component }\tau_{\ge2}g,
 \qquad
 \cofib(\mathcal N^{D}_{E/F})\simeq\fib(\tau_{\ge2}g),
\]
and $g_j$ is split surjective for every $j\ge2$.
If moreover $g_1$ is an isomorphism and $g_2$ is surjective, then $\fib(\tau_{\ge2}g)\simeq Z$, and $\pi_jZ=\ker g_j$ for $j\ge2$.
\end{lemma}
\begin{proof}
Since $(\tau_{\ge2}g)\circ\mathcal N^{D}_{E/F}\circ h^{-1}=hh^{-1}$ is the identity, the map $s=\mathcal N^{D}_{E/F}\circ h^{-1}$ is a section of $\tau_{\ge2}g$.
The section gives the direct-sum decomposition with first component $\tau_{\ge2}g$, and $g_j$ is split surjective for $j\ge2$.
Since $\mathcal N^{D}_{E/F}=s\circ h$ and $h$ is an equivalence, $\cofib(\mathcal N^{D}_{E/F})\simeq\cofib(s)$, and the cofiber of the section $s$ of $\tau_{\ge2}g$ is $\fib(\tau_{\ge2}g)$.
If $g_1$ is an isomorphism and $g_2$ is surjective, the long exact sequence of homotopy groups for $Z=\fib(g)$ gives the vanishing of $\pi_0Z$ and $\pi_1Z$, and the natural map $\fib(\tau_{\ge2}g)\to Z$ is an equivalence.
\end{proof}

We also need splitting fields in which prescribed real places split completely.

\begin{lemma}\label{lem:paper-galois-splitting-real}
Let $F$ be a global field, and let $S$ be a set of real places of $F$ at which $D$ splits.
Then there is a finite Galois splitting field $E/F$ of $D$ in which every place in $S$ splits completely.
\end{lemma}
\begin{proof}
If $S$ is empty, the Galois closure of a separable maximal subfield of $D$ suffices.
Otherwise $F$ is a number field, and we use weak approximation in $D$ to choose an element $\alpha$ that is close, at each place in $S$, to a matrix with distinct real eigenvalues.
The reduced characteristic polynomial of $\alpha$ is separable and has only real roots at each place in $S$.
Since $D$ is a division algebra, $F(\alpha)$ is a field, and it is a maximal subfield of degree $\deg D$ in which every place in $S$ splits completely.
The Galois closure of $F(\alpha)$ has the same property.
\end{proof}

\begin{theorem}\label{thm:paper-global-split}
Let $F$ be a number field and assume that $D$ splits at all real places, or else invert $2$ throughout.
Let $g$ be the comparison of Theorem~\ref{thm:paper-global-comparison}, and put $Z=\fib(g)$.
Let $E/F$ be a finite Galois splitting field of $D$ with group $G$, and if $D$ splits at all real places, assume that every real place of $F$ splits completely in $E$.
Then the composite $(\tau_{\ge2}g)\,\mathcal N^{D}_{E/F}$ is an equivalence.
Consequently,
\[
 \tau_{\ge2}K(D)\simeq\tau_{\ge2}K(F)\oplus Z
 \quad\text{with first component }\tau_{\ge2}g,
 \qquad Z\simeq\cofib(\mathcal N^{D}_{E/F}).
\]
For every $j\ge2$, $g_j$ is split surjective and $\pi_jZ=\ker g_j$.
In particular $Z\in\mathrm{Sp}_{\ge3}$.
\end{theorem}
\begin{proof}
Such a field $E$ exists by Lemma~\ref{lem:paper-galois-splitting-real}, applied to the set of all real places if $D$ splits at all of them and to the empty set otherwise.

We first compare the norms on the completed factors.
For $\ell\mid \deg D$, we restrict to $E$ the rank-$1$ class $x$ used in the comparison.
Since $x_E$ is restricted from $F$, Lemma~\ref{lem:paper-etale-rank-one} gives a coherently $G$-equivariant equivalence
\[
 m_{x_E}\colon L_{K(1),\ell}K(E)\xrightarrow{\ \sim\ }L_{K(1),\ell}K(D_E).
\]
Since $D_E$ is split, Morita equivalence and the Quillen--Lichtenbaum comparison for $E$ identify its $2$-connective cover with an equivariant equivalence
\[
 \alpha_E\colon(\tau_{\ge2}K(E))^\wedge_\ell\xrightarrow{\ \sim\ }(\tau_{\ge2}K(D_E))^\wedge_\ell.
\]
The projection formula of Lemma~\ref{lem:paper-etale-rank-one} gives
\[
   (\tau_{\ge2}g)^\wedge_\ell\,(\mathcal N^{D}_{E/F})^\wedge_\ell\,
           (\alpha_E)_{hG}=(\mathcal N_{E/F})^\wedge_\ell
\]
as maps of $K(F)^\wedge_\ell$-modules in the derived $\ell$-complete category.
Here the homotopy orbits are taken in the derived $\ell$-complete category, and reducing them modulo $\ell$ gives the same result as taking homotopy orbits of the reductions modulo $\ell$ before completion.
This applies at odd primes in the ramified real case as well.

After inverting $\deg D$, multiplication by the restriction to $E$ of $[D]/\deg D\in K_0(D)[1/\deg D]$ is a $G$-equivariant rank-$1$ equivalence $K(E)[1/\deg D]\to K(D_E)[1/\deg D]$.
By the projection formula, the homotopy orbits of this equivalence composed with $\mathcal N^{D}_{E/F}$ and then with $N_{D/F}/\deg D$ give $\mathcal N_{E/F}$.
Proposition~\ref{prop:paper-derived-norm} now proves that $(\tau_{\ge2}g)\mathcal N^{D}_{E/F}$ is an equivalence rationally and modulo all relevant primes, and hence is an equivalence with the stated coefficients.

Finally, $g_1$ and $g_2$ are isomorphisms by Lemma~\ref{lem:paper-global-low-degrees}, so Lemma~\ref{lem:paper-split-from-norm} yields both formulas and $\pi_jZ=\ker g_j$ for $j\ge2$.
Since $g_2$ is injective, $\pi_2Z=0$, and hence $Z\in\mathrm{Sp}_{\ge3}$.
On each completed factor the projection formula also gives the splitting as $K(F)^\wedge_\ell$-modules.
\end{proof}

\subsubsection{Splitting over global function fields by constant field transfer}
Over global function fields, we obtain the corresponding splitting (Theorem~\ref{thm:paper-global-function-split}) by transfer from a constant field extension.

Let $F$ be a global function field with constant field $k=\F_q$.
Let $f_v$ be the residue degree of a place $v$, and choose a finite constant extension $E=Fk'$ such that $s=[k':k]$ is a common multiple of the integers $\deg D\cdot f_v$ for the places $v$ where the local invariant of $D$ is nonzero.
At each such $v$, the local degree $s/\gcd(s,f_v)$ is divisible by $\deg D$, so $D_E$ has zero local invariants and splits by the local--global principle for the Brauer group.

\begin{lemma}\label{lem:paper-global-constant-norm}
For any finite constant extension $E/F$ with group $G=\Gal(E/F)$ and every $j\ge2$,
\[
 H_i(G,K_j(E))=0\quad(i>0),\qquad
 N_{E/F}:K_j(E)_G\xrightarrow{\sim}K_j(F).
\]
\end{lemma}
\begin{proof}
Let $\ell\ne p=\operatorname{char}k$.
We first show $H^2(F,\Q_\ell/\Z_\ell(m))=0$ for $m\ge2$.
Let $C/k$ be the smooth proper geometrically connected model of $F$.
By the Weil weights \cite{Deligne1974}, Frobenius minus one on $H^1(C_{\bar k},\Z_\ell(m))$ is invertible on the rationalized lattice.
Thus it is surjective on $H^1(C_{\bar k},\Q_\ell/\Z_\ell(m))$.
The finite field has cohomological dimension one, so the Hochschild--Serre edge map gives
\[
 H^2(C,\Q_\ell/\Z_\ell(m))\xrightarrow{\sim}H^0(k,H^2(C_{\bar k},\Q_\ell/\Z_\ell(m)))
 =H^0(k,\Q_\ell/\Z_\ell(m-1)).
\]
The Gysin map of a closed point $v$ becomes corestriction
\[
 H^0(k_v,\Q_\ell/\Z_\ell(m-1))\longrightarrow H^0(k,\Q_\ell/\Z_\ell(m-1)).
\]
Indeed, the trace of a geometric point class is one, and summing over conjugate points is corestriction.
For the arithmetic Frobenius $\Phi$ and $f=[k_v:k]$, corestriction is $1+\Phi+\cdots+\Phi^{f-1}$, which on $\Q_\ell/\Z_\ell(m-1)$ is multiplication by a nonzero integer and hence surjective.
The identity $(\Phi-1)(1+\Phi+\cdots+\Phi^{f-1})=\Phi^f-1$ implies surjectivity on invariants as well.
Also $H^1(k_v,\Q_\ell/\Z_\ell(m-1))=0$, so purity and localization give
\[
\begin{aligned}
 H^0(k_v,\Q_\ell/\Z_\ell(m-1))&\longrightarrow H^2(C,\Q_\ell/\Z_\ell(m))\\
 &\longrightarrow H^2(C-\{v\},\Q_\ell/\Z_\ell(m))\longrightarrow H^1(k_v,\Q_\ell/\Z_\ell(m-1)).
\end{aligned}
\]
Thus $H^2(C-\{v\},\Q_\ell/\Z_\ell(m))=0$, and the same argument on smaller open subsets of $C$ gives $H^2(F,\Q_\ell/\Z_\ell(m))=0$ by passage to the direct limit.

We next compute the torsion of $K_j(F)$ for $j\ge2$.
Harder's theorem and localization with closed support give $K_j(F)\otimes\Q=0$ for $j\ge2$ \cite[Corollary 2.3.2]{Harder}.
For an abelian group $M$ and a prime $\ell$, we write $M\{\ell\}$ for the $\ell$-primary torsion subgroup of $M$.
Since $[F:F^p]=p$, Geisser--Levine's calculation with $p$-coefficients shows that $K_j(F)\{p\}$ vanishes in this range \cite{GeisserLevine}.
Since $K_j(F)$ is torsion, the direct limit $K_j(F)\otimes(\Q_\ell/\Z_\ell)$ of the groups $K_j(F)/\ell^r$ vanishes, so the direct limit of the mod-$\ell^r$ Bockstein sequences gives $K_j(F;\Q_\ell/\Z_\ell)\simeq K_{j-1}(F)\{\ell\}$ for $j\ge3$.
We calculate the coefficient groups by the Quillen--Lichtenbaum comparison and \'etale descent \cite{ClausenMathew}.
Over $F\otimes_k\bar k$, which is a field since $k$ is algebraically closed in $F$, the \'etale $K$-theory coefficients are concentrated in even degrees, and the $H^2$ vanishing just proved leaves only $H^0(F,\Q_\ell/\Z_\ell(m))$ in total degree $2m$ and $H^1(F,\Q_\ell/\Z_\ell(m+1))$ in total degree $2m+1$.
Geometric connectedness gives $H^0(F,\Q_\ell/\Z_\ell(m))=H^0(k,\Q_\ell/\Z_\ell(m))$.
Compatibly with transfer, we have
\[
 \begin{aligned}
 K_{2m-1}(F)\{\ell\}&\simeq H^0(k,\Q_\ell/\Z_\ell(m))       &&(m\ge2),\\
 K_{2m}(F)\{\ell\}&\simeq H^1(F,\Q_\ell/\Z_\ell(m+1))      &&(m\ge1).
 \end{aligned}
\]
The same identifications hold over $E$.

Finally we verify the exactness of the periodic resolution for the cyclic group $G$.
In odd degree we set $M=\Q_\ell/\Z_\ell(m)$, and in even degree we set $M=H^1(F\otimes_k\bar k,\Q_\ell/\Z_\ell(m+1))$.
In both degrees, we have
\[
 K_j(F)\{\ell\}\simeq M^\Phi,\qquad K_j(E)\{\ell\}\simeq M^{\Phi^s},
\]
where $j=2m-1$ in odd degree, $j=2m$ in even degree, and $s=[E:F]$.
In even degree this follows from $H^1(k,\Q_\ell/\Z_\ell(m+1))=0$ and the Hochschild--Serre edge map.
Since $F\otimes_k\bar k$ has $\ell$-cohomological dimension at most one, vanishing of $H^2(F,\Q_\ell/\Z_\ell(m+1))$ and $H^2(E,\Q_\ell/\Z_\ell(m+1))$ implies that $\Phi-1$ and $\Phi^s-1$ are surjective on $M$.
In odd degree both are nonzero integer multiplications on $\Q_\ell/\Z_\ell(m)$ and thus surjective.
Put $N_s=1+\Phi+\cdots+\Phi^{s-1}$.
Since $\Phi^s-1$ is surjective and $N_s(\Phi-1)=(\Phi-1)N_s=\Phi^s-1$, the operator $N_s$ is surjective on $M$ and induces a surjection $N_s\colon M^{\Phi^s}\to M^\Phi$.
If $z\in M^{\Phi^s}$ and $N_sz=0$, we write $z=(\Phi-1)y$.
Then $(\Phi^s-1)y=N_sz=0$, so the kernel of $N_s\colon M^{\Phi^s}\to M^\Phi$ is $(\Phi-1)M^{\Phi^s}$.
These two checks give the exactness of the periodic resolution in positive degrees, so higher group homology vanishes and the norm on coinvariants is an isomorphism on $\ell$-primary torsion.
Summing over the primes gives the claim.
\end{proof}

\begin{theorem}\label{thm:paper-global-function-split}
Let $F$ be a global function field.
Let $g$ be the comparison of Theorem~\ref{thm:paper-global-comparison}, and put $Z=\fib(g)$.
Let $E=Fk'$ be a finite constant extension splitting $D$, with group $G=\Gal(E/F)$.
Then the composite $(\tau_{\ge2}g)\,\mathcal N^{D}_{E/F}$ is an equivalence.
Consequently,
\[
 \tau_{\ge2}K(D)\simeq\tau_{\ge2}K(F)\oplus Z
 \quad\text{with first component }\tau_{\ge2}g,
 \qquad Z\simeq\cofib(\mathcal N^{D}_{E/F}).
\]
For every $j\ge1$, $g_j$ is surjective, and for $j\ge2$,
\begin{equation}\label{eq:paper-global-function-quotient}
 \pi_jZ\simeq K_j(D)\big/\theta_E K_j(E)
  =K_j(D)\Big/\sum_{\substack{L/F\ \text{finite separable}\\ D_L\ \text{split}}}
                            \theta_L K_j(L).
\end{equation}
\end{theorem}
\begin{proof}
We first compute the homotopy orbits.
The orbit spectral sequence
\[
 H_i(G,\pi_j\tau_{\ge2}K(D_E))
 \ \Longrightarrow\ \pi_{i+j}\bigl((\tau_{\ge2}K(D_E))_{hG}\bigr)
\]
has a finite filtration in each total degree by connectivity.
Under the usual identification $\pi_jK(D_E)=K_j(E)$ by Morita equivalence, the action of each $\gamma\in G$ agrees with the field action.
Indeed, composing the semilinear action of $\gamma$ with the inverse of the field action gives an $E$-linear automorphism of $D_E\simeq M_{\deg D}(E)$, which is inner by Skolem--Noether and hence acts as the identity on $K$-theory.
By Lemma~\ref{lem:paper-global-constant-norm}, only the row $i=0$ remains, so the edge map identifies $\pi_j\bigl((\tau_{\ge2}K(D_E))_{hG}\bigr)$ with $K_j(E)_G$, and under this identification $\pi_j(\mathcal N^{D}_{E/F})$ is induced by the splitting transfer $\theta_E$.
The transfer identity $g_j\theta_E=N_{E/F}$ therefore shows that $(\tau_{\ge2}g)\mathcal N^{D}_{E/F}$ induces an isomorphism on every homotopy group.
By Lemma~\ref{lem:paper-global-low-degrees}, $g_1$ and $g_2$ are isomorphisms, so Lemma~\ref{lem:paper-split-from-norm} gives the direct-sum decomposition, $\cofib(\mathcal N^{D}_{E/F})\simeq Z$, and the surjectivity of $g_j$ for every $j\ge1$.

Finally we identify $\pi_jZ$.
For any finite separable splitting field $L/F$, the compositum $EL/L$ is a finite constant extension, so $N_{EL/L}\colon K_j(EL)\to K_j(L)$ is surjective for $j\ge2$ by Lemma~\ref{lem:paper-global-constant-norm}.
Transitivity of transfer gives $\theta_LK_j(L)\subseteq\theta_EK_j(E)$, and $E$ itself is a finite separable splitting field.
Together with the long exact sequence of the splitting, this gives~\eqref{eq:paper-global-function-quotient}.
\end{proof}

\subsubsection{The third \texorpdfstring{$K$}{K}-group and real signatures}
This subsubsection computes $K_3(D)$ with integral coefficients (Theorem~\ref{thm:global-k3}).
Together with the splittings above, it shows that the comparison fiber is $3$-connected (Corollary~\ref{cor:global-fibre-connectivity}).

Let $J_F$ be the image of $K_3^M(F)\to K_3(F)$.
For a number field $F$, the signatures at the real places give an isomorphism $J_F\simeq\bigoplus_{v\ \mathrm{real}}\Z/2$, and for a global function field, $J_F=0$ \cite{BassTate,MerkurjevSuslinK3}.
For a number field, write $J_{F}^{\mathrm{ram}}\subset J_F$ for the subgroup corresponding to the summands at the real places where $D$ is ramified, and $J_{F}^{\mathrm{sp}}$ for the subgroup corresponding to the summands at the real places where $D$ splits.
For a global function field, put $J_F^{\mathrm{ram}}=0$.
In any splitting field of $D$, every place above a real place at which $D$ is ramified is complex; hence restriction $K_3(F)\to K_3(L)$ kills $J_F^{\mathrm{ram}}$ for every splitting field $L$, and restriction is defined on $K_3(F)/J_F^{\mathrm{ram}}$.

\begin{theorem}\label{thm:global-k3}
If $F$ is a global function field with constant field $\F_q$, there is a unique isomorphism
\[
 K_3(D)\xrightarrow{\ \sim\ }K_3(F)
\]
compatible with restriction to $K_3(L)$ for every finite separable splitting field $L/F$; moreover, restriction from the constant field is an isomorphism $K_3(\F_q)\simeq K_3(F)$, and $K_3(\F_q)$ is cyclic of order $q^2-1$.
If $F$ is a number field, there is a unique isomorphism
\[
 K_3(D)\xrightarrow{\ \sim\ }K_3(F)/J_{F}^{\mathrm{ram}}
\]
compatible with restriction to $K_3(L)$ for every finite Galois splitting field $L/F$ in which every real place of $F$ at which $D$ splits is completely split.
These isomorphisms are compatible with restriction and transfer for finite extensions.
\end{theorem}
\begin{proof}
\emph{Step 1: reduction to weights two and three.}
By Theorem~\ref{thm:paper-global-weight-one}, $\CH^1(F,3;D)=0$, and the stages $\mathcal{F}^rK(D)$ with $r\ge4$ vanish in degrees at most three (over a global function field, the homotopy coniveau tower is available by Lemma~\ref{lem:app-twisted-tower}), so
\[
 \CH^3(F,3;D)\longrightarrow \pi_3 \mathcal{F}^2K(D)
   \longrightarrow \CH^2(F,3;D)\longrightarrow 0,\qquad
 \pi_3\mathcal{F}^2K(D)\twoheadrightarrow K_3(D).
\]
On the commutative side, $\CH^1(F,4)=0$ as well, so the commutative edge map $\pi_3\mathcal{F}^2K(F)\to K_3(F)$ is an isomorphism and the image of weight three is $J_F$ \cite[Corollary 6.1.4]{KahnLevine2010}.

\emph{Step 2: weight two.}
By Propositions~\ref{prop:global-two-chow-rows} and \ref{prop:global-real-weight-two}, the weight-two reduced norm is an isomorphism integrally.
For a finite separable extension $L/F$, transfer shows that the kernel of restriction on $K_3^{\mathrm{ind}}(F)$ is killed by the extension degree.
The coefficient sequence and $H^0_{\mot}(F,\Q(2))=0$ identify the torsion subgroup of $K_3^{\mathrm{ind}}(F)$ with $H^0(F,\Q/\Z(2))$, and restriction is inclusion of invariants \cite{Voevodsky}; over a global function field, the characteristic-primary part is treated in Step~5.
Hence restriction is injective on $K_3^{\mathrm{ind}}(F)$.
Since $\CH^2(F,3)\simeq K_3^{\mathrm{ind}}(F)$ and the weight-two reduced norm is an isomorphism, restriction of $\CH^2(F,3;D)$ to a splitting field is also injective, with the same image as on the commutative side.

\emph{Step 3: weight three.}
As in the proof of Theorem~\ref{thm:paper-global-weight-one} for $\CH^2(k,2;D)$, pushforward from a maximal subfield at each admissible closed point shows that the group $\CH^3(F,3;D)$ is generated by the transfers of $\CH^3(E,3)=K_3^M(E)$ over finite splitting fields $E/F$ \cite{NesterenkoSuslin}.
For global function fields, these groups $K_3^M$ vanish.
For number fields, we choose a finite Galois splitting field $L/F$ in which every real place of $F$ at which $D$ splits is completely split; it exists by Lemma~\ref{lem:paper-galois-splitting-real}.
Every real place of any finite splitting field $E/F$ lies over a split real place of $F$, so $EL/E$ splits completely at every real place.
Restriction duplicates the real signatures of $K_3^M$ on the real branches, and transfer adds them over the branches.
Thus $K_3^M(EL)\to K_3^M(E)$ is surjective, and transitivity of transfer gives
\[
 K_3^M(L)_{\Gal(L/F)}\twoheadrightarrow \CH^3(F,3;D)
      \xrightarrow{\rho} K_3^M(F).
\]
The composite is the usual norm, and it is an isomorphism onto $J_{F}^{\mathrm{sp}}$: each Galois orbit of real branches corresponds to one split real place, and the unit signature on one branch maps to the unit signature below; here $K_3^M(F)\simeq J_F$ by the real signatures \cite{BassTate}.
Since restriction duplicates signatures, $\CH^3(F,3;D)\to J_L$ is injective with image $J_L^{\Gal(L/F)}$.
The commutative side has kernel $J_{F}^{\mathrm{ram}}$ and the same image \cite[Chapter II, Theorem 2.1(3)]{BassTate}, \cite[Lemmas 6.6.3--6.6.4]{KahnLevine2010}.

\emph{Step 4: assembly.}
For a number field, let $L$ be as in Step~3; for a global function field, let $L$ be any finite separable splitting field.
If an element of $\pi_3\mathcal{F}^2K(D)$ restricts to zero in $K_3(L)$, its weight-two image is zero, so it comes from weight three, where restriction is also injective.
Since this map factors through $K_3(D)$, one obtains $\pi_3\mathcal{F}^2K(D)\simeq K_3(D)$.
Given $x\in K_3(D)$, choose $y\in K_3(F)$ whose image in $K_3^{\mathrm{ind}}(F)=\CH^2(F,3)$ is $\rho$ of the weight-two image of $x$.
For global function fields, $L/F$ is any finite separable splitting field and $J_L=0$, so matching the weight-two images already makes the restrictions of $x$ and $y$ to $K_3(L)$ agree.
For number fields, the difference of the restrictions of $x$ and $y$ to $K_3(L)$ lies in $J_L^{\Gal(L/F)}$, which is the image of weight three on both sides, so after correcting $y$ by an element of $J_F$ the two restrictions agree.
The same argument with the two sides exchanged shows that the images of $K_3(D)$ and $K_3(F)$ in $K_3(L)$ agree.
The kernel on the commutative side is $J_F^{\mathrm{ram}}$, which is zero for a global function field.
Thus $K_3(D)$ and $K_3(F)/J_F^{\mathrm{ram}}$ both inject into $K_3(L)$ with equal images, so there is exactly one isomorphism between them compatible with restriction to $K_3(L)$.

\emph{Step 5: global function fields.}
For global function fields, weight three vanishes.
Here $J_F=0$, so $K_3^{\mathrm{ind}}(F)=K_3(F)$, and its characteristic-primary torsion vanishes by Proposition~\ref{prop:app-frobenius}(2); for the other primary parts we use Beilinson--Lichtenbaum comparison.
Finite-coefficient comparison and Harder's vanishing give, for $\ell\ne\operatorname{char} F$,
\[
 \CH^2(F,3)_{(\ell)}
 = H^1_{\mot}(F,\Z_{(\ell)}(2))
 \simeq H^0(F,\Q_\ell/\Z_\ell(2)).
\]
The action on the right-hand side factors through the constant field, whose Frobenius acts by multiplication by $q^2$.
Hence $K_3(F)$ is cyclic of order $q^2-1$, and restriction from the constant field is an isomorphism \cite{Harder,QuillenFinite}.

\emph{Step 6: independence and functoriality.}
Independence of the splitting field follows by comparing restriction images in a compositum.
Restriction and transfer for arbitrary finite extensions descend to the signature quotients by duplicating and adding signatures at ramified real places.
For a finite extension $E/F$, compatibility with the isomorphisms constructed from restriction images follows from the base-change formula and Morita equivalence over each field component of $E\otimes_F L$.
In particular, $\theta_E$ from a splitting field $E/F$ corresponds to the composite $K_3(E)\xrightarrow{N_{E/F}}K_3(F)\to K_3(F)/J_{F}^{\mathrm{ram}}$ on the commutative side.
\end{proof}

\begin{corollary}\label{cor:global-k3-structure}
Let $F$ be a number field with $r_2$ complex places, let $r_{\mathrm{sp}}$ be the number of real places at which $D$ splits, and put $w_2=|H^0(F,\Q/\Z(2))|$.
Then $K_3(D)$ has free rank $r_2$, and its torsion subgroup is $\Z/w_2$ if $r_{\mathrm{sp}}=0$ and $\Z/(2w_2)\oplus(\Z/2)^{r_{\mathrm{sp}}-1}$ if $r_{\mathrm{sp}}>0$.
\end{corollary}
\begin{proof}
By Theorem~\ref{thm:global-k3}, $K_3(D)\simeq K_3(F)/J_F^{\mathrm{ram}}$, and the free rank of $K_3(F)$ is $r_2$ \cite{Borel}.
Since $J_F^{\mathrm{ram}}$ is finite, the torsion of $K_3(D)$ is $T/J_F^{\mathrm{ram}}$, where $T$ is the torsion of $K_3(F)$.
The group $T$ is an extension of $H^0(F,\Q/\Z(2))=\Z/w_2$ by $J_F$.
For a lift $z\in T$ of a generator (it is torsion since $J_F$ is finite), $w_2z=\{-1,-1,-1\}$ \cite[\S11]{MerkurjevSuslinK3}, the element of $J_F$ that is nonzero at every real place.
Hence $T/J_F^{\mathrm{ram}}$ is an extension of $\Z/w_2$ by $J_F^{\mathrm{sp}}$ in which $w_2\bar z$ is the vector with all components one; it is zero if there are no split real places and of order two otherwise, which gives the stated torsion.
\end{proof}

\begin{corollary}\label{cor:global-fibre-connectivity}
The fiber $Z$ of the comparison map $g$ belongs to $\mathrm{Sp}_{\ge4}$.
If $D$ is ramified at a real place, the same holds for the fiber of the comparison map with $2$ inverted.
\end{corollary}
\begin{proof}
By Theorems~\ref{thm:paper-global-split} and~\ref{thm:paper-global-function-split}, $g_3$ is surjective.
By Theorem~\ref{thm:global-k3}, its source and target are abstractly isomorphic over $\Z$, or over $\Z[1/2]$ under real ramification.
Both are finitely generated: for number fields, localization makes $K_3(F)$ an image of $K_3(\mathcal O_F)$, since $K_2$ of finite fields vanishes, and for global function fields $K_3(F)$ is finite by the same theorem.
A surjection between isomorphic finitely generated modules over these Noetherian rings is an isomorphism, so $g_3$ is an isomorphism.
By Lemma~\ref{lem:paper-global-low-degrees}, $g_1$ and $g_2$ are isomorphisms.
The splitting in degrees at least two therefore makes the homotopy groups of $Z$ vanish in degrees at most three.
\end{proof}

\begin{remark}
At each prime $\ell\ne\operatorname{char}F$ to which Proposition~\ref{prop:global-two-chow-rows} applies, the comparison in degree three also agrees with the weight-two reduced norm.
Indeed, the rank-$1$ trivialization acts as the identity on the geometric weight-two layer of \'etale $K$-theory, and the reduced norm on $\CH^2(F,3;D)_{(\ell)}$ is the comparison on that layer.
Let $L$ be a splitting field as in Theorem~\ref{thm:global-k3}.
If $\ell\mid\deg D$, restriction carries the completed comparison $m_x^{-1}$ to $m_{x_L}^{-1}$; the latter acts as the identity in positive degrees after split Morita identification by Lemma~\ref{lem:paper-global-product-vanishing}.
If $\ell\nmid\deg D$, flat base change gives the same compatibility because $N_{D_L/L}$ is $\deg D$ times the split Morita comparison.
Since $K_3(L)_{(\ell)}$ injects into its derived completion, the two restrictions agree in $K_3(L)_{(\ell)}$.
At these primes $J_F{}_{(\ell)}=0$, so the restriction description in Theorem~\ref{thm:global-k3} identifies $g_3$ with the weight-two reduced norm.
\end{remark}

When $D$ is ramified at a real place, the comparison of Theorem~\ref{thm:paper-global-comparison} is defined after inverting $2$.
Integrally, $K_1(D)$ is the subgroup of $F^\times$ consisting of the elements positive at every ramified real place.
The group $K_1(D)$ has no two-torsion corresponding to the element $-1$ of $K_1(F)$.
\begin{example}\label{ex:paper-global-k3-real-ramification}
Let $F=\Q$, $D=(-1,-1)$, and $E=\Q(i)$.
There is no integral map $K_3(D)\to K_3(\Q)$ carrying all splitting transfers to the corresponding field norms.
\end{example}
\begin{proof}
By Theorem~\ref{thm:global-k3}, $K_3(D)=\Z/24$ and $K_3(\Q)=\Z/48$, and $\theta_E$ corresponds to the composite $K_3(E)\xrightarrow{N_{E/\Q}}K_3(\Q)\to K_3(\Q)/J_{\Q}^{\mathrm{ram}}$.
Lift a generator of $K_3^{\mathrm{ind}}(\Q)=\Z/24$ to $z\in K_3(\Q)=\Z/48$ and put $y=\res_{E/\Q}z$.
The surjection $K_3(\Q)=\Z/48\to K_3^{\mathrm{ind}}(\Q)=\Z/24$ has kernel $J_{\Q}=\{0,24\}$, so $z$ generates $\Z/48$.
Then $N_{E/\Q}y=2z$ has order $24$, while its image $\theta_E(y)$ in the signature quotient $K_3(D)$ has order $12$.
Since $12\theta_E(y)=0$ but $12N_{E/\Q}y\ne0$, no homomorphism $K_3(D)\to K_3(\Q)$ carries $\theta_E(y)$ to $N_{E/\Q}y$ \cite[\S11]{MerkurjevSuslinK3}.
\end{proof}

\subsubsection{The basic comparison in prime index}
For a prime $\ell\ne\operatorname{char}F$ at which the $\ell$-primary index of $D$ is $\ell$, the comparison can be analyzed directly by comparing twisted and commutative cycle complexes.
This case, treated in Proposition~\ref{prop:global-squarefree-spectrum}, is an input for the degree bounds in the next subsubsection.

\begin{proposition}\label{prop:global-squarefree-spectrum}
For a prime $\ell\mid \deg D$ with $\ell\ne\operatorname{char}F$, and $\ell\ne2$ if $D$ is ramified at a real place, such that $v_\ell(\deg D)=1$, the map $g$ of Theorem~\ref{thm:paper-global-comparison} is an equivalence after $\ell$-localization.
Consequently, if the prime-to-characteristic part of the index is squarefree, $g$ is an equivalence in the coefficient range of that theorem.
\end{proposition}

\begin{proof}
Since the rational comparison holds, the integral comparison after $\ell$-localization reduces to the $\ell$-completed comparison.
We split the primary components of the Brauer class other than the $\ell$-primary component over a separable extension of degree prime to $\ell$.
After $\ell$-localization, restriction to this extension is split injective, since its composite with transfer is multiplication by the degree, and the same holds for the cofiber of the localization unit $K(D)\to L_{K(1),\ell}K(D)$.
Since $m_x$ is an equivalence, it suffices to show that this localization unit is an equivalence after completion in positive degrees, and by the split injectivity this may be checked after restriction.
We are therefore reduced to the case where the division algebra in the Brauer class has degree $\ell$.
In positive characteristic, both composites of restriction and transfer for a finite purely inseparable extension are multiplication by its degree, which is prime to $\ell$.
Passing to the direct limit, we therefore move the calculation modulo $\ell^r$ to the perfect closure, as required by \cite[Theorem 8.2.2]{KahnLevine2010}, which is used below.
The \'etale side is unchanged by purely inseparable extension.

We first compare the twisted and commutative cycle complexes in positive weight.
By \cite[Theorem 8.2.2]{KahnLevine2010}, the reduced cycle map in positive weight $q$ is an isomorphism in motivic degrees $s<q$.
In degree $q$ it is injective, with cokernel the image of
\[
 K_q^M(F)\xrightarrow{[D]\cup-}
 H^{q+2}_{\et}(F,\Z/\ell(q+1)).
\]
Since $q+2\ge3$, the target $H^{q+2}_{\et}(F,\Z/\ell(q+1))$ vanishes when the $\ell$-cohomological dimension is at most two.
For $\ell=2$ over a number field with real places, restriction of the target to all real closures is injective, and $[D]=0$ at each real closure, so the image is again zero.
In degrees $s>q$, both twisted and commutative terms vanish because there are no cycles in simplicial degree $2q-s<q$.
Thus the positive-weight map is a quasi-isomorphism of complexes over a field.
It remains a quasi-isomorphism modulo $\ell^r$ and commutes with restriction to real closures \cite[Theorem 6.5.3]{KahnLevine2010}, \cite{Jannsen,Voevodsky}.

We compare the homotopy coniveau tower modulo $\ell^r$ with the tower of global sections of its \'etale hypersheafification.
The layer comparison is the map from the twisted motivic complex over a field to \'etale cohomology with coefficients $\Z/\ell^r(q)$, placed in homotopy degree $2q-s$.
This follows from the Kahn--Levine identification of the layers and from the fact that \'etale sheafification of the reduced cycle map identifies the coefficient sheaves.
By the quasi-isomorphism just proved, the comparison cofiber for $q>0$ identifies with the commutative one.
By Beilinson--Lichtenbaum comparison and vanishing of motivic cohomology in degrees $>q$, its cohomology occurs only for $s>q$, where it equals $H^s_{\et}(F,\Z/\ell^r(q))$.
For $q\ge2$, these cohomology groups vanish in cohomological dimension two.
In the real-relative case, restriction to the real closures is also an isomorphism in degrees $s\ge3$, so the real-relative comparison cofiber vanishes.
For $q=1$, both the absolute and real-relative cofibers lie in homotopy degrees $2-s\le0$ \cite[Theorem 6.5.5, \S6.9]{KahnLevine2010}.

Convergence of the towers is controlled by the following degree estimate.
The term with supports in codimension at least $q$ has no supports in simplicial degree below $q$, so the $q$th stage of the homotopy coniveau tower is concentrated in degrees at least $q$, and so are its geometric stalks.
The stalk of a layer modulo $\ell^r$ is a finite-coefficient Tate sheaf in degree $2q$.
In any fixed degree, the geometric homotopy group is computed using sufficiently many but finitely many layers, and is therefore a finite $\ell$-primary coefficient group.
Consequently, after \'etale hypersheafification, taking global sections or real-relative sections lowers connectivity by at most the cohomological-dimension bound two.
The tails at sufficiently high stages do not contribute to a fixed total degree.
Layerwise comparison therefore shows that $\pi_j$ of the comparison cofiber for the entire tower vanishes for $j\ge2$ and that $\pi_1$ of this cofiber injects into $\pi_1$ of the weight-zero comparison cofiber.

It remains to check weight zero.
The algebraic layer places the rank lattice $\deg D\cdot\Z$ with coefficients modulo $\ell^r$ in degree zero, and the comparison on $\pi_0$ is
\[
 \deg D\cdot\Z/\ell^r(\deg D\cdot\Z)\longrightarrow\Z/\ell^r,\qquad [\deg D]\longmapsto [\deg D].
\]
The \'etale layer has no positive homotopy groups, so $\pi_1$ of the weight-zero comparison cofiber is the kernel of this map.
In the real-relative case, $D$ splits over each real closure, and all local weight-zero comparison cofibers lie in degrees at most zero.
The same kernel occurs in $\pi_1$ of the relative fiber.
Under the coefficient transition, this kernel at level $r+v_\ell(\deg D)$ maps to zero at level $r$.
Thus the inverse limit of $\pi_1$ of the entire comparison cofiber is zero.
Since $\pi_2$ of the comparison cofiber vanishes at every finite stage, the Milnor exact sequence also makes $\pi_1$ of the derived-completed comparison cofiber zero.
This proves that the completed comparison is an isomorphism in every positive degree.
In the real-relative case, this proves the comparison for the relative terms.
Over the real closures $D$ is split, so Morita equivalence and the Quillen--Lichtenbaum comparison for the real closures give the comparison for the real terms.
If the relative and real terms are isomorphic in positive degrees, the long exact sequence gives isomorphisms of the absolute terms in degrees at least two; in degree one we use that $g_1=\Nrd_1$ is an isomorphism (Lemma~\ref{lem:paper-global-low-degrees}).
Applying the arithmetic square shows that $g$ is an equivalence after $\ell$-localization.
For the second assertion, the map $g[1/\deg D]=N_{D/F}/\deg D$ is an equivalence because $i_DN_{D/F}=N_{D/F}i_D=(\deg D)^2$.
If $p=\operatorname{char}F$ divides $\deg D$, Proposition~\ref{prop:app-frobenius} gives $K_i(D;\Z/p)=K_i(F;\Z/p)=0$ for $i\ge2$, so the $p$-completed comparison is the completion of $g_1=\Nrd_1$, and $g_1$ is an isomorphism by Lemma~\ref{lem:paper-global-low-degrees}.
Together with the first assertion at the other primes dividing $\deg D$, the arithmetic square for the primes dividing $\deg D$ shows that $g$ is an equivalence.
\end{proof}

\subsubsection{Prime decomposition and degree bounds for the comparison fiber}
We describe the primary factors of the comparison fiber using the localization fibers $B_\ell(D)$ (Theorem~\ref{thm:paper-global-fiber-bounds}).

For $\ell\mid \deg D$ with $\ell\ne\operatorname{char}F$, write
\[
 e_\ell=\ell^{v_\ell(\deg D)},\qquad
 B_\ell(D)=\fib\!\left(K(D)_{(\ell)}\longrightarrow L_{E(1),\ell}K(D)\right).
\]

\begin{theorem}\label{thm:paper-global-fiber-bounds}
Let $F$ be a global field; if $D$ is ramified at a real place, invert $2$ throughout and exclude $\ell=2$.
Let $g\colon\tau_{\ge1}K(D)\to\tau_{\ge1}K(F)$ be the comparison of Theorem~\ref{thm:paper-global-comparison}, and put $Z=\fib(g)$.
Then the fiber decomposes as
\begin{equation}\label{eq:paper-global-fiber-primes}
 Z\simeq
 \prod_{\substack{\ell\mid \deg D,\ \ell\ne\operatorname{char}F\\
                  \ell\ne2\ \text{if $D$ is ramified at a real place}}}
       \tau_{\ge2}B_\ell(D).
\end{equation}
For each factor,
\begin{equation}\label{eq:paper-global-fiber-amplitude}
 e_\ell\,\pi_j B_\ell(D)=0\quad (j\ge1),
 \qquad
 B_\ell(D)\in\mathrm{Sp}_{\le e_\ell-3}.
\end{equation}
The factor $\tau_{\ge2}B_\ell(D)$ is zero if $e_\ell=\ell$.
\end{theorem}

\begin{proof}
The functor $L_{E(1),\ell}$ is the smashing localization combining rationalization and height one.
Since $\ell$ is invertible over $F$, the $\mathrm{THH}(F)$-module structure makes multiplication by $\ell$ an equivalence on $\mathrm{THH}$ of algebras and schemes over $F$.
It is an equivalence of cyclotomic spectra and hence also on $\mathrm{TC}$, so $\mathrm{TC}_{(\ell)}$ is rational.
In the definition of Selmer $K$-theory, the map from $\mathrm{TC}_{(\ell)}$ to its $L_{E(1),\ell}$ localization is therefore an equivalence, which gives an equivalence $K^{\mathrm{Sel}}_{(\ell)}\simeq L_{E(1),\ell}K$ compatible with the comparison maps \cite[Definition 6.1]{ClausenMathew}.

We first bound the exponent of the factors.
We choose a finite separable splitting field $E/F$ of degree $\deg D$.
The virtual $\ell$-cohomological dimension of $E$ is at most two, so \cite[Theorem 6.18]{ClausenMathew} gives $B_\ell(E)\in\mathrm{Sp}_{\le-2}$.
By split Morita equivalence, we have $B_\ell(D_E)\simeq B_\ell(E)$.
We have $N_{D_E/D}\res_{E/F}=\deg D$ on both $K(D)$ and $L_{E(1),\ell}K(D)$.
Thus $\deg D$ kills $\pi_jB_\ell(D)$ for $j\ge1$.
These groups are $\ell$-local, so their exponent divides $e_\ell$.

For the upper degree bound, we split the prime-to-$\ell$ Brauer component over a finite separable extension $F'/F$ of degree prime to $\ell$.
This does not change the $\ell$-primary index, since $\ind(D_{F'})\mid\deg D\mid[F':F]\,\ind(D_{F'})$.
Restriction and transfer divided by the extension degree exhibit $\pi_jB_\ell(D)$ as a direct summand of $\pi_jB_\ell(D_{F'})$.
Thus by Morita equivalence it suffices to treat a division algebra of degree $e_\ell$.
For $X=\SB(D)$, we have the Quillen decomposition \cite[Section 8]{QuillenHigher}
\[
 K(X)\simeq\bigoplus_{i=0}^{e_\ell-1}K(D^{\otimes i}).
\]
Applying this decomposition to the localization unit and its fiber, we realize $B_\ell(D)$ as a summand of $B_\ell(X)$.
The residue fields of points of $X$ have virtual $\ell$-cohomological dimension at most $2+\dim X=e_\ell+1$.
The same theorem of Clausen and Mathew gives $B_\ell(X)\in\mathrm{Sp}_{\le e_\ell-3}$ and hence \eqref{eq:paper-global-fiber-amplitude}.

Next we take the derived $\ell$-completion of $g$.
The construction in Theorem~\ref{thm:paper-global-comparison} uses $L_{K(1),\ell}K(F)\simeq L_{K(1),\ell}K(D)$, so the completed comparison fiber in degrees at least two is
\[
 \tau_{\ge2}\fib\!\left(
   (\tau_{\ge1}K(D))^\wedge_\ell
      \longrightarrow \tau_{\ge1}L_{K(1),\ell}K(D)
 \right).
\]
Replacing $\tau_{\ge1}K(D)$ by $K(D)$ leaves this unchanged, since the completion of $\tau_{\le0}K(D)$ is concentrated in degrees at most one by $(-)^\wedge_\ell=\varprojlim_r(-)/\ell^r$ and the Moore coefficient sequence.
Removing $\tau_{\ge1}$ in the target leaves the same $2$-connective fiber.
The height-one arithmetic square gives $(L_{E(1),\ell}K(D))^\wedge_\ell\simeq L_{K(1),\ell}K(D)$, so the completed comparison fiber in this range is $\tau_{\ge2}B_\ell(D)^\wedge_\ell$.

The bounded exponent above implies that $\tau_{\ge1}B_\ell(D)$ is derived $\ell$-complete.
The Bockstein torsion terms have transition maps given by multiplication by $\ell$ and form a pro-zero system.
The quotient terms stabilize in each degree, so the Milnor sequence recovers the homotopy groups after completion.
Since the completion of $\tau_{\le0}B_\ell(D)$ is concentrated in degrees at most one, $\tau_{\ge2}B_\ell(D)^\wedge_\ell\simeq\tau_{\ge2}B_\ell(D)$.
On the other hand, $g$ is an isomorphism in degrees one and two by Lemma~\ref{lem:paper-global-low-degrees}, so $Z$ is $2$-connective and its completion is too.
We obtain $Z^\wedge_\ell\simeq\tau_{\ge2}B_\ell(D)$.

The construction of $g[1/\deg D]$ gives $Z[1/\deg D]=0$.
In characteristic $p$, the completed positive-degree comparison is an equivalence at $p$ by the rank-$1$ Frobenius comparison of Proposition~\ref{prop:app-frobenius} and the degree-one reduced norm.
The arithmetic square of $Z$ therefore leaves only the completed factors in \eqref{eq:paper-global-fiber-primes}.
When $e_\ell=\ell$, restrict to $F'$ as above and pass by Morita equivalence to the division algebra $D'$ in the Brauer class of $D_{F'}$, which has degree $\ell$.
Proposition~\ref{prop:global-squarefree-spectrum}, applied to $D'$ over $F'$, shows that the comparison for $D'$ is an equivalence after $\ell$-localization.
Its completed comparison fiber in degrees at least two is $\tau_{\ge2}B_\ell(D')$, as identified above, so $\tau_{\ge2}B_\ell(D_{F'})\simeq\tau_{\ge2}B_\ell(D')=0$.
Restriction and transfer, whose composite is an $\ell$-local unit, show that $\tau_{\ge2}B_\ell(D)$ is a summand of this zero spectrum.
\end{proof}

\begin{corollary}\label{cor:paper-global-fiber-finite}
Let $F$ be a number field, and let $\ell$ be a prime dividing $\deg D$, with $\ell\ne2$ if $D$ is ramified at a real place.
Then every homotopy group of $\tau_{\ge2}B_\ell(D)$ is a finite $\ell$-group.
\end{corollary}
\begin{proof}
We choose a maximal $\mathcal O_F$-order $\Lambda\subset D$ and a nonempty open subscheme $U\subset\Spec\mathcal O_F$ such that $\mathcal A=\Lambda|_U$ is Azumaya and $U$ contains no place above $\ell$.
The groups $K_j(\Lambda)$ are finitely generated \cite{Kuku1982}.
Maximal orders are hereditary and the support terms at the deleted places are $K$-theories of finite residue division rings, so localization gives finite generation of $K_j(\mathcal A)$ as well \cite[Section 5]{QuillenHigher}.
The support term for $K(\mathcal A)\to K(D)$ is the direct sum of $K(k_v)$ over the closed points $v$ of $U$.
Each $B_\ell(k_v)$ is concentrated in degrees at most $-2$, and height-one localization commutes with direct sums and localization sequences \cite{QuillenFinite,ClausenMathew}.
Thus $B_\ell(\mathcal A)\to B_\ell(D)$ is an isomorphism in degrees at least two.
The equivalence $Z^\wedge_\ell\simeq\tau_{\ge2}B_\ell(D)$ obtained in the proof of Theorem~\ref{thm:paper-global-fiber-bounds} (and $Z_{(\ell)}\to Z^\wedge_\ell$ is an equivalence, since the other factors of the finite product have bounded exponent prime to $\ell$) commutes, after $\ell$-completion, with the canonical maps to $(\tau_{\ge2}K(D))^\wedge_\ell$, the common source of the completed comparison and of the completed localization unit.
For $j\ge2$, the map $\pi_jB_\ell(D)\to\pi_jB_\ell(D)^\wedge_\ell$ is an isomorphism by the bounded exponent in \eqref{eq:paper-global-fiber-amplitude}.
Since the completed comparison splits by Theorem~\ref{thm:paper-global-split}, the composite $\pi_jB_\ell(D)\to K_j(D)_{(\ell)}\to\pi_jK(D)^\wedge_\ell$ is injective for $j\ge2$, and hence so is $\pi_jB_\ell(D)\to K_j(D)_{(\ell)}$.
The latter map factors through $\pi_jB_\ell(\mathcal A)\to K_j(\mathcal A)_{(\ell)}$, exhibiting $\pi_jB_\ell(D)$ as a subgroup of a finitely generated $\Z_{(\ell)}$-module.
By \eqref{eq:paper-global-fiber-amplitude}, the group $\pi_jB_\ell(D)$ is finite.
\end{proof}

\begin{corollary}\label{cor:paper-global-fiber-prime-four}
If $e_\ell=\ell$ or $e_\ell=4$ for every $\ell\mid \deg D$ with $\ell\ne\operatorname{char}F$, then $g$ is an equivalence (with the coefficient convention of Theorem~\ref{thm:paper-global-fiber-bounds}).
Over a global function field, we have $\pi_jZ=0$ for $j<4$, and the $\ell$-primary part of $\pi_jZ$ can survive only in the range
\[
 4\le j\le e_\ell-3.
\]
\end{corollary}
\begin{proof}
If $e_\ell=4$ then $e_\ell-3=1$ and $\tau_{\ge2}B_\ell(D)=0$.
The case $e_\ell=\ell$ is the last assertion of Theorem~\ref{thm:paper-global-fiber-bounds}.
Over a global function field, $\pi_jZ=0$ for $j<4$ by Corollary~\ref{cor:global-fibre-connectivity}.
We combine this with \eqref{eq:paper-global-fiber-amplitude}.
\end{proof}

\begin{remark}
Equations~\eqref{eq:paper-global-function-quotient} and \eqref{eq:paper-global-fiber-primes} describe the same fiber by transfers and by height-one localization, respectively.
The latter gives the prime decomposition and an exponent bound for $K_j(D)/\theta_EK_j(E)$ over a global function field.
\end{remark}

\subsubsection{Invariance of the fiber under powers of the Brauer class}
We use Adams operations and the prime decomposition~\eqref{eq:paper-global-fiber-primes} to compare the fibers for different powers of the Brauer class.

\begin{proposition}
\label{prop:paper-brauer-power-fiber}
Let $F$ be a global field and let $D$ be a central division algebra over $F$.
Let $a\ge1$ be an integer with $\gcd(a,\per D)=1$.
Let $g_D$ and $g_{D^{\otimes a}}$ be the comparisons of Theorem~\ref{thm:paper-global-comparison} for $D$ and for the division algebra in the Brauer class of $D^{\otimes a}$, and put $Z_D=\fib(g_D)$ and $Z_{D^{\otimes a}}=\fib(g_{D^{\otimes a}})$.
For number fields with real ramification, the comparison fibers are taken with $2$ inverted, as in that theorem.
There is a noncanonical equivalence of spectra
\[
 Z_D\simeq Z_{D^{\otimes a}}.
\]
Consequently, two central simple algebras whose classes generate the same cyclic subgroup of $\Br(F)$ have equivalent comparison fibers under the same coefficient convention.
Moreover, if nonnegative integers $i,j$ satisfy $\gcd(i,\per D)=\gcd(j,\per D)$, then $Z_{D^{\otimes i}}\simeq Z_{D^{\otimes j}}$, where $D^{\otimes0}=F$.
\end{proposition}
The proof of Proposition~\ref{prop:paper-brauer-power-fiber} is given after Lemmas~\ref{lem:paper-brauer-power-correspondence}--\ref{lem:paper-brauer-power-descent} below.

For the auxiliary comparison, let $F$ be any field and let $D$ be a central division algebra over $F$ with $\deg D\ge2$, and take $a\ge2$.
Write $D^i=D^{\otimes i}$ and $X=\SB(D^{\mathrm{op}})$.
We regard the universal right $D^{\mathrm{op}}$-ideal as a left $D$-bundle, and let $\phi_i:K(D^i)\to K(X)$ be the maps defined by its tensor powers.
The Quillen decomposition is
\[
 \phi=\bigoplus_{i=0}^{\deg D-1}\phi_i:
 P=\bigoplus_{i=0}^{\deg D-1}K(D^i)\xrightarrow{\sim}K(X).
\]
The negative universal bundles give a semiorthogonal decomposition ordered by $i=\deg D-1,\ldots,0$; reorder the $K$-theory direct sum in increasing order.
We take the inclusions and projections to be exact functors on the differential graded enhancements of the categories of perfect complexes \cite[\S8]{QuillenHigher}, \cite[Theorem 4.1]{Bernardara2009}.
The same universal bundle defines $\phi_i$ for all $i\ge0$, giving
\[
 \phi^{(a)}=\bigoplus_{i=0}^{\deg D-1}\phi_{ai}:
 P_a=\bigoplus_{i=0}^{\deg D-1}K(D^{ai})\longrightarrow K(X).
\]

\begin{lemma}\label{lem:paper-brauer-power-correspondence}
If $\gcd(a,\per D)=1$, then $\phi^{(a)}[1/a]\colon P_a[1/a]\to K(X)[1/a]$ is an equivalence.
\end{lemma}
\begin{proof}
Since $\per D\mid\deg D$ and multiplication by $a$ permutes $\Z/(\per D)$, the classes $[D^i]$ and $[D^{ai}]$ for $0\le i<\deg D$ agree as multisets, so permuting components and applying Morita equivalences gives $P_a\simeq P$.
Composing the projections in the Quillen decomposition with $\phi_{ai}$, we obtain derived tensor products with perfect bimodule complexes between the components.
We use the $K_0$-classes of these bimodule complexes as correspondences, with composition given by derived tensor product.
Taking Morita ranks after splitting gives a faithful representation of the resulting endomorphism ring of correspondences as
\[
 \Pi=\left\{(n_{ij})\in M_{\deg D}(\Z):
 n_{ij}\in\ind\bigl((D^i)^{\mathrm{op}}\otimes_F D^j\bigr)\Z\right\}
\]
with ordinary matrix multiplication as composition.
By additivity of $K$-theory, equalities of correspondences yield equalities of homotopy classes of maps of spectra \cite[Proposition 2.25(i), \S10, (10.1)]{TabuadaVanDenBerghSeparable}.

After splitting, put $t=[\mathcal O(-1)]$.
The maps $\phi$ and $\phi^{(a)}$ compare the two lists $1,t,\ldots,t^{\deg D-1}$ and $1,t^a,\ldots,t^{a(\deg D-1)}$ in $\Z[t]/(1-t)^{\deg D}$.
Writing $h=1-t$, we have $1-t^a=ah+O(h^2)$, so the degree-$r$ factor in the $h$-filtration is $a^r$.
The matrix comparing the two lists therefore has determinant
\[
                    \pm a^{\deg D\,(\deg D-1)/2}.
\]
Because $\gcd(a,\per D)=1$, the algebras $D^i$ and $D^{ai}$ have the same splitting fields and hence the same index.
Through the chosen permutation and the Morita equivalences $P_a\simeq P$, the map $\phi^{-1}\phi^{(a)}$ becomes an endomorphism of $P$ whose matrix in $\Pi$ is the split $K_0$-matrix above up to a permutation of components, so its determinant is $\pm a^{\deg D(\deg D-1)/2}$.
The Cayley--Hamilton identity expresses its inverse as a polynomial in it with coefficients in $\Z[1/a]$, so the inverse belongs to $\Pi[1/a]$.
Realizing this inverse in $K$-theory shows that $\phi^{(a)}[1/a]$ is an equivalence.
\end{proof}

\begin{lemma}\label{lem:paper-brauer-power-adams}
The stable Adams operation $\psi_X^a$ is an equivalence on $K(X)[1/a]$.
\end{lemma}
\begin{proof}
Let $KGL$ be the motivic $\PP^1$-spectrum representing algebraic $K$-theory, let $f_q$ be the $q$-effective cover, and let $s_q=\cofib(f_{q+1}\to f_q)$ be the $q$th slice.
Write $\operatorname{Fil}^qK(X)=(f_qKGL)(X)$; in particular, $\operatorname{Fil}^0K(X)\simeq K(X)$.
The unit $\eta:\mathbf1[1/a]\to KGL[1/a]$ induces an isomorphism on the zeroth slice, and $\psi^a\eta=\eta$ implies $s_0\psi^a=1$.
Since $\psi^a(b)=ab$ for the Bott element $b$, we have $s_r\psi^a=a^r$ for every $r$.
Put $m=\dim X$.
The $q$th stage of the motivic filtration $\operatorname{Fil}^\bullet K(X)[1/a]$ belongs to $\mathrm{Sp}_{\ge q-m}$.
A fixed $\pi_n$ can be computed in the finite quotient $\operatorname{Fil}^0/\operatorname{Fil}^q$ for $q>n+m$, and $a^r$ is a unit on each of its slices.
Thus $\psi_X^a$ induces an isomorphism on every $\pi_n$.
For the stable operation and its action on the Bott element, see \cite[Example 9.6]{AnnalaHoyoisIwasa}; the unit equivalence on slices and the convergence estimate are \cite[Corollaries 9.5 and 9.7]{BachmannElmantoMorrow}.
Integer localization commutes with effective covers and slices because the effective subcategories have compact generators.
\end{proof}

\begin{lemma}\label{lem:paper-brauer-power-descent}
For every $i\ge0$ there are Adams operations $\psi_i^a\colon K_n(D^i)\to K_n(D^{ai})$, $n\ge0$, such that
\[
 (\psi_X^a)_*(\phi_i)_*=(\phi_{ai})_*\psi_i^a\qquad\text{on }K_n(D^i)[1/a].
\]
\end{lemma}
\begin{proof}
We choose a finite Galois splitting field $L/F$, put $G=\Gal(L/F)$, and choose $\gamma$-semilinear isomorphisms $\xi_\gamma$ for $\gamma\in G$ realizing the descent of $D_L=\operatorname{End}_L(W)$, with $\xi_1=1$.
Write $\xi_\gamma\xi_\delta=c_{\gamma,\delta}\xi_{\gamma\delta}$.
The crossed product of the resulting normalized $2$-cocycle $c$ acts on $W$ and has commutant $D$, so the Brauer class of $c$ is $-[D]$.
Let $\mathcal M_i$ be the exact category of finite-dimensional $L$-vector spaces $V$ equipped with $\gamma$-semilinear isomorphisms $\zeta_\gamma:V\to V$, for $\gamma\in G$, satisfying
\[
                     \zeta_\gamma\zeta_\delta=c_{\gamma,\delta}^{i}\zeta_{\gamma\delta}.
\]
Morphisms are $L$-linear maps commuting with the $\zeta_\gamma$.
The category $\mathcal M_i$ is the category of left modules over the crossed product of $c^i$.
The opposite algebra of this crossed product has Brauer class $i[D]$, so $\mathcal M_i$ is Morita equivalent to the category of right $D^i$-modules.
Tensor products add the exponents of the multipliers, and $a$th exterior powers of flags multiply them by $a$.
After splitting, a right $D^i$-module has the form $V\otimes_L(W^*)^{\otimes i}$.
Since $(W^*)^{\otimes i}$ has multiplier $c^{-i}$, the multiplier of $V$ is $c^i$.
The variety $X_L$ parametrizes lines in $W^*$, so $\mathcal O(-1)$ has multiplier $c^{-1}$, which cancels in $V\otimes_L\mathcal O(-i)$.
This ordinary descent defines $\phi_i$.
The $\zeta_\gamma$ are compatible with the natural isomorphisms for tensor products and exterior powers of flags, so the natural isomorphisms descend to natural isomorphisms on the original exact categories.
This system of graded exact categories gives exterior-power operations $\lambda_i^a$ and Adams operations $\psi_i^a:K_n(D^i)\to K_n(D^{ai})$ on the groups \cite[Definition 1.1, Definition 2.1, Proposition 2.2, (2.8), Theorem 4.1]{KockZanchettaPower}.
The Newton polynomial has total tensor degree $a$, and in positive degrees $\psi_i^a=(-1)^{a-1}a\lambda_i^a$; the target of $\psi_i^a$ is therefore $K_n(D^{ai})$.
Applying the functorial construction of exterior-power operations to the natural isomorphisms above gives the displayed identity for the Adams operation on $K_n(X)$ defined by exterior powers.

It remains to identify this Adams operation with the stable operation $\psi_X^a$ on $K_n(X)[1/a]$.
Precomposing the stable operation with $KGL\to KGL[1/a]$, and composing the operation defined by the Newton polynomial in exterior powers with the map to $[1/a]$-coefficients, we obtain two pointed maps $\Z\times\mathrm{Gr}\to\Omega^\infty KGL[1/a]$.
Both send the class of a line bundle to its $a$th power, so by the splitting principle they induce the same natural transformation on $K_0$ of smooth $F$-schemes.
The proof of \cite[Lemma 1.2.10]{Riou2010} that restriction on $K_1$ of finite Grassmannians is surjective remains valid after inverting $a$, so the target satisfies the condition of \cite[Theorem 1.1.6]{Riou2010}, and the two maps agree.
Hence they agree on every $K_n\to K_n[1/a]$, and by the universal property of localization of coefficients they agree on $K_n(X)[1/a]$ \cite[Example 9.6]{AnnalaHoyoisIwasa}.
\end{proof}

\begin{proof}[Proof of Proposition~\ref{prop:paper-brauer-power-fiber}]
The assertion is immediate if $a=1$ or $\deg D=1$, so we assume that $a\ge2$ and $\deg D\ge2$, and we use the notation above.
By Lemmas~\ref{lem:paper-brauer-power-correspondence}, \ref{lem:paper-brauer-power-adams} and~\ref{lem:paper-brauer-power-descent}, the equivalence of spectra
\[
 (\phi^{(a)})^{-1}\psi_X^a\phi:P[1/a]\xrightarrow{\sim}P_a[1/a]
\]
induces, on all homotopy groups, a diagonal matrix preserving the index $i$.
Each diagonal component is an isomorphism of groups, so the map of spectra defined using inclusion and projection,
\[
 \Psi_D^{(a)}=\operatorname{pr}_1(\phi^{(a)})^{-1}\psi_X^a\phi\operatorname{in}_1:
 K(D)[1/a]\longrightarrow K(D^{\otimes a})[1/a]
\]
is an equivalence.

Now assume that $F$ is a global field.
If $\ell\mid \deg D$ and $\ell\ne\operatorname{char}F$, then $\ell\nmid a$, since $\per D$ and $\deg D$ have the same prime divisors.
We localize the equivalence above at $\ell$ and use naturality of the localization unit for $L_{E(1),\ell}$ to obtain $B_\ell(D)\simeq B_\ell(D^{\otimes a})$.
Unit powers preserve the index, so combining the finitely many factors of~\eqref{eq:paper-global-fiber-primes} gives $Z_D\simeq Z_{D^{\otimes a}}$ without inverting $a$.
If there is real ramification, invert $2$ as in Theorem~\ref{thm:paper-global-comparison}.

Since generators of the same cyclic subgroup of $\Br(F)$ are related by a unit power, the first consequence follows.
If $\gcd(i,\per D)=\gcd(j,\per D)$, then $i[D]$ and $j[D]$ generate the same cyclic subgroup, so the second consequence also follows.
\end{proof}

\subsubsection{Examples with vanishing comparison fiber at arbitrary index}
Example~\ref{ex:paper-global-fiber-rational} gives algebras of arbitrary index whose comparison fiber vanishes.

\begin{example}\label{ex:paper-global-fiber-rational}
Let $k$ be a finite field and put $F=k(t)$.
Let $k_n/k$ be the extension of degree $n$, and let $\sigma$ be its Frobenius.
Then
\[
 D=(k_n(t)/k(t),\sigma,t)
\]
has index $n$, and the comparison $g$ in positive degrees is an equivalence.
\end{example}
\begin{proof}
The local invariant at $t=0$ has order $n$, so the cyclic algebra $D$ has index $n$.
Put $\mathcal A=k_n[u,u^{-1};\sigma]$ and identify its center with $k[t,t^{-1}]$ by $t=u^n$.
After base change to $k_n$, the coefficient algebra decomposes into diagonal idempotents, and $u$ becomes a cyclic permutation matrix whose full cycle multiplies by $t$.
These matrices generate the matrix algebra, so $\mathcal A$ is Azumaya over $k[t,t^{-1}]$.

By the twisted Bass--Heller--Swan decomposition \cite[Corollary 5.3]{KirsteinKremer} and the vanishing of the Nil terms \cite[Corollary 4.2]{KirsteinKremer}, which follows from regularity of the coefficient field $k_n$, we have
\[
 K(\mathcal A)\simeq\cofib\!\left(1-\sigma_*:K(k_n)\longrightarrow K(k_n)\right).
\]
Height-one localization is exact, and taking fibers is an exact functor in a stable category, so the same decomposition holds for $B_\ell$.
The bound $B_\ell(k_n)\in\mathrm{Sp}_{\le-2}$ for finite fields \cite[Theorem 6.18]{ClausenMathew} gives $B_\ell(\mathcal A)\in\mathrm{Sp}_{\le-1}$.
In the localization sequence for $K(\mathcal A)\to K(D)$, the support term is the direct sum of $K(k_v)$ over the closed points $v$ of $\G_{m,k}$.
Indeed, the Azumaya algebra splits over each finite residue field $k_v$, so Morita equivalence and d\'evissage apply \cite[\S5]{QuillenHigher}.
Since $L_{E(1),\ell}$ commutes with direct sums, applying $B_\ell$ to this localization sequence gives a fiber sequence whose support term is the direct sum of the spectra $B_\ell(k_v)$.
Each $B_\ell(k_v)$ is in degrees at most $-2$, so $B_\ell(D)$ is in degrees at most $-1$.
All factors in~\eqref{eq:paper-global-fiber-primes} therefore vanish.
\end{proof}

\begin{remark}
Quillen's finite-field calculation gives $K_{2m-1}(\F_q)=\Z/(q^m-1)$ \cite{QuillenFinite}, so every connective cover of the $K$-theory of a finite field is unbounded above.
Moreover, every nonzero upper truncation $\tau_{\le N}\tau_{\ge1}K(\F_q)$ has a nonzero homotopy group in degree at most three.
By Theorem~\ref{thm:paper-global-fiber-bounds}, the comparison fiber is bounded above, and by Corollary~\ref{cor:global-fibre-connectivity}, it lies in degrees at least four for a global function field.
It remains open to describe the homotopy groups and Postnikov invariants of the comparison fiber in this finite range of degrees for global function fields.
\end{remark}

\subsubsection{Eilenberg--Mac Lane decompositions over number fields}\label{subsec:paper-global-fiber-gem}
We prove the odd-primary part of Theorem~\ref{thm:paper-global-fiber-gem} here and use the Bott element for the $2$-primary part in the next two subsubsections.

A finite generalized Eilenberg--Mac Lane spectrum (finite GEM) is a finite direct sum of suspensions of Eilenberg--Mac Lane spectra.
Let $F$ be a number field, let $D$ be a central division algebra over $F$, and let $Z$ be the fiber of the comparison in Theorem~\ref{thm:paper-global-comparison}.
If there is real ramification, take the comparison and $Z$ with $2$ inverted.

\begin{theorem}\label{thm:paper-global-fiber-gem}
For each prime $\ell\mid\deg D$, put $e_\ell=\ell^{v_\ell(\deg D)}$.
There is a noncanonical equivalence of spectra
\[
 Z \simeq
 \bigoplus_{\ell\mid\deg D}\ \bigoplus_{j=4}^{e_\ell-3}
           \Sigma^j H(\pi_jZ_{(\ell)}).
\]
If there is real ramification, the sum omits $\ell=2$; an empty degree range contributes zero.
Each $\pi_jZ_{(\ell)}$ is a finite $\ell$-group and, by the splitting in Theorem~\ref{thm:paper-global-split}, is the kernel of $K_j(D)_{(\ell)}\to K_j(F)_{(\ell)}$.
\end{theorem}

We prove Theorem~\ref{thm:paper-global-fiber-gem} after the following lemmas.
Fix a prime $\ell$ as in the theorem and put $R=K(F)_{(\ell)}$ and $\mathsf{X}=\tau_{\ge2}B_\ell(D)$.

\begin{lemma}\label{lem:paper-gem-positive-action}
There is a finite layer $E$ of the cyclotomic $\Z_\ell$-extension of $F$, with $G=\Gal(E/F)$, such that the following hold.
\begin{enumerate}
\item The action map $\tau_{\ge2}R\otimes_R\mathsf{X}\to\mathsf{X}$ is nullhomotopic as a map of $R$-modules.
\item The map
\[
 (\tau_{\ge1}K(E)_{(\ell)})_{hG}\otimes_R\mathsf{X}\longrightarrow\mathsf{X},
 \qquad y\otimes z\longmapsto\mathcal N_{E/F}(y)\cdot z,
\]
is nullhomotopic as a map of $R$-modules.
\end{enumerate}
\end{lemma}
\begin{proof}
By Theorem~\ref{thm:paper-global-fiber-bounds}, $\mathsf X\simeq Z_{(\ell)}$ as spectra.
We use the $R$-module structure on $\mathsf X$ induced by the localization fiber defining $B_\ell(D)$.
The localization map is a module map, and restriction, transfer, and their projection formula also induce maps and identities on fibers.

We take $F_\infty$ to be the compositum of $F$ with the cyclotomic $\Z_\ell$-extension of $\Q$, so every real place splits completely.
We take a sufficiently high finite layer $E$ in the cyclotomic $\Z_\ell$-extension $F_\infty/F$, and put $G=\Gal(E/F)$.
The decomposition group of every finite place is open in $\Gal(F_\infty/F)$.
Indeed, away from $\ell$, the Frobenius of the place acts through its absolute norm, whose image in the $\Z_\ell$-quotient of $\Z_\ell^\times$ is nonzero; above $\ell$, the ramification indices in the local cyclotomic extension are unbounded.
We may therefore choose $E$ so that, at each of the finitely many places where the $\ell$-primary Brauer class of $D$ is nontrivial, the local degree is divisible by the order of the local invariant of that class.
The invariants of this primary component vanish at real places, so Brauer--Hasse--Noether implies that $\ind(D_E)$ is prime to $\ell$.
Applying the argument that proves the first relation of \eqref{eq:paper-global-fiber-amplitude} to $D_E$ over $E$, with a separable splitting field of degree $\ind(D_E)$, we obtain $\ind(D_E)\,\pi_jB_\ell(D_E)=0$ for $j\ge1$.
Since $\ind(D_E)$ is prime to $\ell$ and $B_\ell(D_E)$ is $\ell$-local, we have $\tau_{\ge2}B_\ell(D_E)=0$.

Since $R$ is connective and $\mathsf X=\tau_{\ge2}B_\ell(D)$, by adjunction for connective covers of $R$-modules, a map of $R$-modules from a $2$-connective $R$-module to $\mathsf X$ is nullhomotopic once its composite with $\mathsf X\to B_\ell(D)$ is.
Restriction $\mathsf{X}\to B_\ell(D_E)$ is nullhomotopic as a map of $G$-equivariant $R$-modules, since $\mathsf{X}$ lies in degrees at least $4$ and $B_\ell(D_E)$ in degrees at most $1$, so the $t$-structure on equivariant modules makes the mapping space contractible.
By the projection formula, the map
\[
 (\tau_{\ge2}K(E)_{(\ell)})_{hG}\otimes_R \mathsf{X}\ \longrightarrow\ B_\ell(D),
 \qquad y\otimes z\ \longmapsto\
       N_{D_E/D}(y\cdot\res_{E/F}z)
\]
is nullhomotopic and equals the map $y\otimes z\mapsto\mathcal N_{E/F}(y)\cdot z$ given by the action through the derived norm.
Using the equivalence $(\tau_{\ge2}K(E)_{(\ell)})_{hG}\simeq\tau_{\ge2}R$ of Proposition~\ref{prop:paper-derived-norm} and factorization of the connective tensor product through the truncation, we see that the composite of the action map in~(1) with $\mathsf X\to B_\ell(D)$ is nullhomotopic.
The source of the action map in~(1) is $6$-connective, since $\mathsf X$ lies in degrees at least $4$, so the principle above gives~(1).
The same projection formula holds with $\tau_{\ge2}K(E)_{(\ell)}$ replaced by $\tau_{\ge1}K(E)_{(\ell)}$, so the composite of the map in~(2) with $\mathsf{X}\to B_\ell(D)$ is nullhomotopic as a map of $R$-modules.
The source of the map in~(2) is $5$-connective, since $\mathsf{X}$ lies in degrees at least $4$.
The same principle gives~(2).
\end{proof}

\begin{lemma}\label{lem:paper-gem-retract}
A retract of a finite GEM is a finite GEM.
Moreover, for every prime $\ell$, a bounded-above and bounded-below $H\Z_{(\ell)}$-module splits as a module into the direct sum of the suspended Eilenberg--Mac Lane spectra of its homotopy groups.
\end{lemma}
\begin{proof}
Order the terms by increasing degree and combine terms of the same degree.
Maps from higher to lower degrees vanish by the $t$-structure, so the ideal of off-diagonal terms in the endomorphism ring is nilpotent.
Let $\varepsilon$ be the idempotent defining the retract and $\varepsilon_0$ its diagonal part.
Then $u=\varepsilon\varepsilon_0+(1-\varepsilon)(1-\varepsilon_0)$ is congruent to $1$ modulo the off-diagonal ideal, hence invertible, and satisfies $\varepsilon u=u\varepsilon_0$.
Thus $\varepsilon=u\varepsilon_0u^{-1}$ is conjugate to $\varepsilon_0$.
The diagonal idempotents specify direct summands of the individual coefficient groups.
For the second assertion, $\Z_{(\ell)}$ has global dimension $1$, so the Ext groups of degree at least $2$ containing the Postnikov extensions vanish.
\end{proof}

\begin{proof}[The odd-primary part of Theorem~\ref{thm:paper-global-fiber-gem}]
Tensor $\tau_{\ge2}R\to R\to\tau_{\le1}R$ with $\mathsf{X}$.
By Lemma~\ref{lem:paper-gem-positive-action}(1), the map $\mathsf{X}\to \tau_{\le1}R\otimes_R \mathsf{X}$ is a split monomorphism of $R$-modules.
For odd $\ell$, one has $\tau_{\le1}\mathbb S_{(\ell)}\simeq H\Z_{(\ell)}$.
Truncating the unit map gives $H\Z_{(\ell)}\to\tau_{\le1}R$, so $\tau_{\le1}R\otimes_R \mathsf{X}$ is an $H\Z_{(\ell)}$-module.
We truncate the entire retract diagram to the finite range of degrees supporting the homotopy groups of $\mathsf{X}$.
Lemma~\ref{lem:paper-gem-retract} shows that $\mathsf{X}$ is also a finite GEM.
\end{proof}

\subsubsection{The Bott element and the \texorpdfstring{$2$}{2}-primary decomposition}
The $2$-primary part rests on two lemmas: Lemma~\ref{lem:paper-gem-bott} on the Bott element and Lemma~\ref{lem:paper-gem-mod-two} on the decomposition modulo two.

\begin{lemma}\label{lem:paper-gem-bott}
Let $\ell=2$.
Let $\widetilde\eta\in\pi_2(\mathbb S/2)$ be an element whose Bockstein is the Hopf element $\eta$, and let $b\in K_2(F;\Z/2)$ be its image under the unit map.
The maps $\Sigma^2 \mathsf{X}\to \mathsf{X}/2$ induced by $b$ and $\Sigma \mathsf{X}\to \mathsf{X}$ induced by $\eta$ are nullhomotopic as maps of $R$-modules.
\end{lemma}
\begin{proof}
For the finite cyclotomic layer $E/F$ of Lemma~\ref{lem:paper-gem-positive-action}, $-1\in N_{E/F}E^\times$.
Indeed, at each finite place the local reciprocity map $F_v^\times\to\Gal(F_\infty/F)$ kills finite-order elements, since its target is a torsion-free subgroup of $\Z_2$.
The real places split completely.
Thus $-1$ is a local norm at every place and is a global norm by the Hasse norm theorem for cyclic extensions \cite[\S1]{HasseNorm1931}.

With integral coefficients, put $Y=(\tau_{\ge1}K(E))_{hG}$.
The derived norm in degrees at least $2$ makes $\pi_2Y\to K_2(F)$ surjective.
In degree $1$, $\pi_1Y=(E^\times)_G$, and cyclic Hilbert 90 makes the field norm $\pi_1Y\to F^\times$ injective.
The preimage of $-1$ therefore lies in $\pi_1Y[2]$.
Under the sphere unit, modeled by permutations of finite sets, the transposition representing $\eta$ becomes a matrix of determinant $-1$, so $\eta$ maps to $-1\in K_1(F)$ and the Bockstein of $b$ is $-1$.
We lift the preimage of $-1$ through the Moore coefficient sequence
\[
 0\longrightarrow \pi_2Y/2\longrightarrow\pi_2(Y/2)
   \longrightarrow\pi_1Y[2]\longrightarrow0.
\]
The norm of this lift differs from $b$ by a class in $K_2(F)/2$, and surjectivity of $\pi_2Y\to K_2(F)$ lets us correct the lift, so $b$ lifts to $\pi_2(Y/2)$.
We apply Lemma~\ref{lem:paper-gem-positive-action}(2) to $Y_{(2)}$ and reduce modulo $2$.
It follows that the action of $b$ is zero.
Composing with the Bockstein shows that the action of $\eta$ is zero as well.
\end{proof}

\begin{lemma}\label{lem:paper-gem-mod-two}
The spectrum $\mathsf{X}/2$ is a finite direct sum of suspensions of $H\F_2$.
\end{lemma}
\begin{proof}
Put $\mathsf{W}=\tau_{\le1}R\otimes_R(\mathsf{X}/2)$.
Reducing the nullhomotopy of Lemma~\ref{lem:paper-gem-positive-action}(1) modulo $2$ makes $\mathsf{X}/2$ a retract of $\mathsf{W}$.
If $\iota$ denotes the inclusion of the bottom cell of the Moore spectrum and $\partial$ the projection to its top cell, then $2\operatorname{id}_{\mathbb S/2}=\iota\eta\partial$ \cite[Proposition 4]{SchwedeMoore}.
Indeed, a self-map whose composite with $\iota$ is zero comes from the order-two group $\pi_1(\mathbb S/2)$, and its nonzero element is $\iota\eta\partial$.
The inequality $2\operatorname{id}_{\mathbb S/2}\ne0$ follows from the action of $Sq^2$ on the bottom degree of $\mathbb S/2\wedge\mathbb S/2$ and the Cartan formula.
Lemma~\ref{lem:paper-gem-bott} therefore gives $2\operatorname{id}_{\mathsf{X}/2}=0$.

We regard $\mathsf{W}$ as a $ko_{(2)}$-module through the ring maps
\[
 ko_{(2)}\longrightarrow \tau_{\le1}ko_{(2)}
    \simeq \tau_{\le1}\mathbb S_{(2)}\longrightarrow\tau_{\le1}R.
\]
The middle equivalence comes from the isomorphisms induced by the sphere unit on $\pi_0,\pi_1$.
The action of $\eta$ on $\mathsf{W}$ is obtained from its action on $\mathsf{X}$ in Lemma~\ref{lem:paper-gem-bott} by base change to $\tau_{\le1}R$ and reduction modulo $2$, so it is zero.
Wood's sequence $\Sigma ko\xrightarrow{\eta}ko\to ku$ makes $\mathsf{W}$ a retract of $\mathsf{V}=ku_{(2)}\otimes_{ko_{(2)}}\mathsf{W}$.
Here the connective version of Wood's sequence follows from the two-cell description of $\Sigma^{-2}\mathbb{CP}^2$ in the periodic version and comparison with complexification on connective covers.
In degrees $0,1$, use that $\pi_0ko\to\pi_0ku$ is an isomorphism and $\eta:\pi_0ko\to\pi_1ko$ is surjective \cite[Theorem 9.7]{MathewNaumannNoel2017}.

The image $b_0\in\pi_2(ko/2)$ of $\widetilde\eta$ maps under complexification to the nonzero element of $\pi_2(ku/2)=\F_2$, namely the reduction of the complex Bott element $\beta\in\pi_2ku$.
Indeed, in Wood's sequence reduced modulo $2$, the map $\eta\colon\pi_0(ko/2)\to\pi_1(ko/2)$ is an isomorphism, so $\pi_2(ko/2)\to\pi_2(ku/2)$ is surjective with kernel $\eta\pi_1(ko/2)$, generated by the reduction of the integral class $\eta^2$.
The Bockstein vanishes on this kernel but takes $b_0$ to $\eta\ne0$, so $b_0$ does not lie in the kernel.

Since $b_0$ comes from the sphere, its action $\Sigma^2\mathsf{W}\to\mathsf{W}/2$ is obtained by base change and reduction modulo $2$ from the action in Lemma~\ref{lem:paper-gem-bott}.
After base change to $ku_{(2)}$, the composite $\Sigma^2\mathsf{V}\xrightarrow{\beta}\mathsf{V}\to \mathsf{V}/2$ is therefore zero.
Since $2\operatorname{id}_{\mathsf{V}}=0$ (the nullhomotopy of $2\operatorname{id}_{\mathsf X/2}$ is $R$-linear, so base change gives a $ku_{(2)}$-linear one), the map $\mathsf{V}\to \mathsf{V}/2$ is a split monomorphism of $ku_{(2)}$-modules, so $\beta_{\mathsf{V}}=0$.
As $ku_{(2)}/\beta\simeq H\Z_{(2)}$, the spectrum $\mathsf{V}$ is a retract of $H\Z_{(2)}\otimes_{ku_{(2)}}\mathsf{V}$.
Consequently, $\mathsf{X}/2$ is a retract of an $H\Z_{(2)}$-module.
We truncate the entire retract diagram to the bounded degree range of $\mathsf{X}/2$.
Lemma~\ref{lem:paper-gem-retract} shows that $\mathsf{X}/2$ is a finite GEM.
Finally, $2\operatorname{id}_{\mathsf{X}/2}=0$ implies that every coefficient group is an $\F_2$-vector space.
\end{proof}

\subsubsection{From finite-coefficient to integral decompositions}
It remains to pass from decompositions with finite coefficients to integral ones, which completes the $2$-primary part of Theorem~\ref{thm:paper-global-fiber-gem}.

\begin{lemma}\label{lem:paper-gem-operations}
Let $M,M'$ be finite $\ell$-groups.
If $n\ge2$, then $\ell[HM,\Sigma^nHM']=0$.
\end{lemma}
\begin{proof}
Richter's account of Kochman's calculation shows that $H_q(H\Z;\Z_{(\ell)})$ is annihilated by $\ell$ for $q>0$ \cite[\S6.1, p.~297]{Richter2006}.
The cofiber of the unit $\mathbb S\to H\Z$ lies in degrees at least $2$, and smashing with $H\Z_{(\ell)}$ preserves this connectivity.
Thus $H_0(H\Z;\Z_{(\ell)})=\Z_{(\ell)}$ and $H_1(H\Z;\Z_{(\ell)})=0$.

We may assume $M=\Z/\ell^a$.
By Lemma~\ref{lem:paper-gem-retract}, each finite-degree truncation of $H\Z_{(\ell)}\wedge H\Z$ splits into suspended Eilenberg--Mac Lane spectra.
In the cofiber of multiplication by $\ell^a$, the free summand in degree zero gives $H(\Z/\ell^a)$, and each summand $\Sigma^qHA$ in positive degree gives $\Sigma^qHA\oplus\Sigma^{q+1}HA$, since $\ell A=0$.
Thus $H_0(HM;\Z_{(\ell)})=\Z/\ell^a$, $H_1(HM;\Z_{(\ell)})=0$, and $H_q(HM;\Z_{(\ell)})$ is annihilated by $\ell$ for $q\ge2$; in particular, no extension in positive degrees produces $\Z/\ell^2$.
We compute $[HM,\Sigma^n HM']$ by the extension--restriction adjunction for $H\Z_{(\ell)}$-modules.
Its summands are
\[
 \operatorname{Hom}_{\Z_{(\ell)}}(H_n(HM;\Z_{(\ell)}),M'),
 \qquad
 \operatorname{Ext}^1_{\Z_{(\ell)}}(H_{n-1}(HM;\Z_{(\ell)}),M').
\]
Here we again use the splitting afforded by global dimension $1$.
For $n\ge2$, both summands are annihilated by $\ell$.
\end{proof}

\begin{lemma}\label{lem:paper-gem-detection}
Let $T$ be a bounded-above and bounded-below spectrum such that every $\pi_iT$ is a finite $\ell$-group.
If $T/\ell$ is a direct sum of suspensions of $H\F_\ell$, then $T$ is a finite GEM.
\end{lemma}
\begin{proof}
If a map from a direct sum of suspensions of $H\F_\ell$ is surjective on $\pi_*$, then its target has the same form and the map admits a section.
Indeed, lift a homogeneous basis of the target to the source and realize it by maps from factors of the same degree; this gives an equivalence onto the target.

Put $T_m=\tau_{\le m}T$.
The map $T/\ell\to T_m/\ell$ is surjective on every homotopy group: it is an isomorphism in degrees at most $m$, and in degree $m+1$ the Bockstein gives a surjection onto $\pi_mT[\ell]$.
Thus $T_m/\ell$ is also such a direct sum.

Inductively assume that $T_{m-1}$ is a finite GEM.
Put $M'=\pi_mT$ and write the Postnikov invariant as $k:T_{m-1}\to\Sigma^{m+1}HM'$.
The composites of the maps from $T/\ell$ show that $T_m/\ell\to T_{m-1}/\ell$ is likewise surjective on all homotopy groups.
By the observation above it admits a section, so $k/\ell=0$.
By naturality of reduction modulo $\ell$, the composite of $k$ with the reduction $\Sigma^{m+1}HM'\to\Sigma^{m+1}HM'/\ell$ factors through $k/\ell=0$, so the coefficient cofiber sequence makes $k$ a multiple of $\ell$.
The degree difference from each summand of $T_{m-1}$ to $\Sigma^{m+1}HM'$ is at least $2$, so Lemma~\ref{lem:paper-gem-operations} gives $k=0$.
\end{proof}

\begin{proof}[Completion of the proof of Theorem~\ref{thm:paper-global-fiber-gem}]
For each prime in the decomposition of Theorem~\ref{thm:paper-global-fiber-bounds}, the groups $\pi_jZ_{(\ell)}$ are finite $\ell$-groups supported in degrees $4\le j\le e_\ell-3$ by \eqref{eq:paper-global-fiber-amplitude} and Corollaries~\ref{cor:global-fibre-connectivity} and~\ref{cor:paper-global-fiber-finite}.
When the $2$-primary factor occurs, Lemmas~\ref{lem:paper-gem-mod-two} and~\ref{lem:paper-gem-detection} applied to $T=\mathsf X\simeq Z_{(2)}$ give the $2$-primary part.
\end{proof}

\subsubsection{The first remaining group}
By Corollary~\ref{cor:global-fibre-connectivity}, the first homotopy group of the comparison fiber that can be nonzero is $\pi_4Z$.
Over a number field, we describe $\pi_4Z_{(\ell)}$ in Proposition~\ref{prop:global-fibre-pi4} when $\ell$ is odd or $F$ has no real places.

\begin{proposition}\label{prop:global-fibre-pi4}
Let $F$ be a number field, let $\ell$ be a prime, and assume either that $\ell$ is odd or that $F$ has no real places.
There is a natural short exact sequence
\[
\begin{split}
 0\longrightarrow&
 \coker\bigl[\CH^1(F,5;D)_{(\ell)}
      \xrightarrow{d_2}\CH^2(F,4;D)_{(\ell)}\bigr]\\
 \longrightarrow&\,\pi_4Z_{(\ell)}
 \longrightarrow\CH^1(F,4;D)_{(\ell)}
 \longrightarrow 0.
\end{split}
\]
Here $d_2$ is the differential in the twisted motivic spectral sequence.
\end{proposition}
\begin{proof}
All groups in this proof are localized at $\ell$.
We index the spectral sequence of the homotopy coniveau tower $\mathcal{F}^rK(D)$ by weight $r$ and total degree $m$; then $d_s$ goes from $(r,m)$ to $(r+s-1,m-1)$.
Weight four in total degree four vanishes by Proposition~\ref{prop:global-two-chow-rows} and the Bass--Tate calculation.
Weight zero and weights at least five also vanish.

First compare weight three with the actual comparison map $g_4$.
Since $\pi_4\mathcal{F}^4K(D)=\pi_3\mathcal{F}^4K(D)=0$, one has
\[
 \pi_4 \mathcal{F}^3K(D)\simeq\CH^3(F,4;D).
\]
On the commutative side, $H^0_{\mot}(F,\Z_{(\ell)}(2))$ and $H^{-1}_{\mot}(F,\Z_{(\ell)}(2))$ vanish by vanishing of the rational groups and Beilinson--Lichtenbaum comparison with divisible coefficients.
Thus weight two vanishes in total degrees four and five.
Weight one also vanishes in these total degrees, so the commutative edge map $\CH^3(F,4)\to K_4(F)$ is an isomorphism.

With finite coefficients, Morita equivalence over $F^{\mathrm{sep}}$ and the motivic--\'etale comparison identify the geometric weight-$r$ layer of the \'etale sheafified homotopy coniveau tower $\mathcal{F}^rK(D)$ with $\Z/\ell^a(r)[2r]$.
The layer comparison is given by reduced rank.
Since these layers are concentrated in degree $2r$ and the connectivity of the tails of the support tower tends to infinity, the sheafified $r$th stage is the geometric $2r$-connective cover.
Multiplication by the rank-$1$ class therefore preserves the stages and acts as the identity on the weight-three layer.
Since $\cd_\ell F\le2$, the stages of weight at least four are $8$-connective at geometric points, $6$-connective on sections and $5$-connective after the derived inverse limit, so their $\pi_4$ and $\pi_3$ vanish as needed in the long exact sequence.
Consequently, the target of the weight-three comparison on $\pi_4$ is $H^2_{\mathrm{cont}}(F,\Z_\ell(3))$.
Thus, after derived completion, the composite $\pi_4\mathcal{F}^3K(D)\to K_4(D)\xrightarrow{g_4}K_4(F)$ equals the composite
\[
 \pi_4\mathcal{F}^3K(D)\simeq\CH^3(F,4;D)\longrightarrow\CH^3(F,4)\longrightarrow K_4(F),
\]
where $\CH^3(F,4;D)\to\CH^3(F,4)$ is the weight-three reduced norm and $\CH^3(F,4)\to K_4(F)$ is the commutative edge map \cite[\S6.9]{KahnLevine2010}, \cite{ClausenMathew}.

For an $\ell$-local spectrum, the map to its derived completion has rational fiber, so its kernel on homotopy groups is divisible.
By Lemma~\ref{lem:paper-global-no-divisible}, $K_4(F)_{(\ell)}$ therefore injects into $\pi_4(K(F)^\wedge_\ell)$, and the comparison above holds with $\Z_{(\ell)}$-coefficients.
Write $\mathcal{F}^3K_4(D)$ for the image of $\pi_4\mathcal{F}^3K(D)$ in $K_4(D)$.
Proposition~\ref{prop:global-two-chow-rows} with $r=3$ implies that $g_4$ maps $\mathcal{F}^3K_4(D)$ isomorphically onto $K_4(F)$.
In particular, all of weight three survives in $K_4(D)$, and $\pi_4Z_{(\ell)}\simeq K_4(D)_{(\ell)}/\mathcal{F}^3K_4(D)_{(\ell)}$.

Next we consider the differential from weight one,
\[
 d_2:\CH^1(F,4;D)\longrightarrow\CH^2(F,3;D).
\]
The comparison map from its source goes to the \'etale weight-one term $H^{-2}_{\et}(F,\Z/\ell^a(1))$, which vanishes because it is in negative cohomological degree.
The comparison on the target is injective: the weight-two reduced norm is an isomorphism by Proposition~\ref{prop:global-two-chow-rows}.
The commutative group $\CH^2(F,3)_{(\ell)}$ is finitely generated over $\Z_{(\ell)}$, so it is detected by its quotients modulo $\ell^a$, and by the coefficient sequence and Beilinson--Lichtenbaum comparison each quotient $\CH^2(F,3)_{(\ell)}/\ell^a$ injects into $H^1_{\et}(F,\Z/\ell^a(2))$.
Since the comparison with the \'etale spectral sequence commutes with $d_2$, the differential is zero.
The target of the next differential, $\CH^3(F,3;D)$, vanishes by Proposition~\ref{prop:global-two-chow-rows} with $r=3$ and \cite[Chapter II, Theorem 2.1(3)]{BassTate}, and higher targets vanish by the dimension condition on supports.
Differentials entering weight one also vanish, leaving $\CH^1(F,4;D)$ in weight one.
Differentials out of weight two vanish for the same reason, and the only incoming differential is $d_2$ from total degree five in weight one.
The resulting two-step filtration in weights one and two gives the displayed short exact sequence \cite[Corollary 6.1.4, \S6.9]{KahnLevine2010}.
\end{proof}

\begin{problem}\label{prob:global-fibre-vanishing}
Do all the finite groups $\pi_jZ_{(\ell)}$ in Theorem~\ref{thm:paper-global-fiber-gem} vanish over number fields?
A nonzero group requires $e_\ell\ge7$, since otherwise the range $4\le j\le e_\ell-3$ is empty, and $e_\ell\ne\ell$, since Theorem~\ref{thm:paper-global-fiber-bounds} gives $\tau_{\ge2}B_\ell(D)=0$ for $e_\ell=\ell$; so the smallest possible values of $e_\ell$ are $8=2^3$ and $9=3^2$.
In degree four, at odd primes or over a field without real places, Proposition~\ref{prop:global-fibre-pi4} reduces this question to conditions on the twisted higher Chow groups in that proposition.
\end{problem}
